% Options for packages loaded elsewhere
\PassOptionsToPackage{unicode}{hyperref}
\PassOptionsToPackage{hyphens}{url}
\documentclass[
]{article}
\usepackage{xcolor}
\usepackage[margin=0.85in]{geometry}
\usepackage{amsmath,amssymb}
\setcounter{secnumdepth}{5}
\usepackage{iftex}
\ifPDFTeX
  \usepackage[T1]{fontenc}
  \usepackage[utf8]{inputenc}
  \usepackage{textcomp} % provide euro and other symbols
\else % if luatex or xetex
  \usepackage{unicode-math} % this also loads fontspec
  \defaultfontfeatures{Scale=MatchLowercase}
  \defaultfontfeatures[\rmfamily]{Ligatures=TeX,Scale=1}
\fi
\usepackage{lmodern}
\ifPDFTeX\else
  % xetex/luatex font selection
\fi
% Use upquote if available, for straight quotes in verbatim environments
\IfFileExists{upquote.sty}{\usepackage{upquote}}{}
\IfFileExists{microtype.sty}{% use microtype if available
  \usepackage[]{microtype}
  \UseMicrotypeSet[protrusion]{basicmath} % disable protrusion for tt fonts
}{}
\makeatletter
\@ifundefined{KOMAClassName}{% if non-KOMA class
  \IfFileExists{parskip.sty}{%
    \usepackage{parskip}
  }{% else
    \setlength{\parindent}{0pt}
    \setlength{\parskip}{6pt plus 2pt minus 1pt}}
}{% if KOMA class
  \KOMAoptions{parskip=half}}
\makeatother
\usepackage{longtable,booktabs,array}
\usepackage{caption}
\captionsetup[table]{skip=6pt}
\newcounter{none} % for unnumbered tables
\usepackage{calc} % for calculating minipage widths
% Correct order of tables after \paragraph or \subparagraph
\usepackage{etoolbox}
\makeatletter
\patchcmd\longtable{\par}{\if@noskipsec\mbox{}\fi\par}{}{}
\makeatother
\setlength{\emergencystretch}{3em} % prevent overfull lines
\providecommand{\tightlist}{%
  \setlength{\itemsep}{0pt}\setlength{\parskip}{0pt}}
\usepackage[round,comma,sort&compress]{natbib}
\bibliographystyle{plainnat}
% ============================================================================
% REQUIRED PACKAGES - Essential for document rendering
% ============================================================================

% Mathematical typesetting (required for equations and symbols)
\usepackage{amsmath,amssymb}          % Mathematical symbols and environments
\usepackage{amsfonts}                 % Additional math fonts
\usepackage{amsthm}                   % Theorem environments

% Graphics and page layout (required for figures and formatting)
\usepackage{graphicx}                 % Include graphics (REQUIRED for figures)
\usepackage{geometry}                 % Page margins (set via pandoc --variable)
\usepackage{float}                    % Better float placement

% Tables (required for table formatting)
\usepackage{booktabs}                 % Professional tables
\usepackage{longtable}                % Long tables spanning pages
\usepackage{array}                    % Advanced table formatting

% PDF features (required for cross-references and metadata)
\usepackage{url}                      % URL formatting
\usepackage{hyperref}                 % Hyperlinks and cross-references
% natbib is loaded once by the Pandoc template; options are applied in _render_pdf_override._postprocess_tex

% ============================================================================
% PACKAGES - Improve formatting and functionality
% ============================================================================

% Table enhancements (optional but recommended)
\usepackage{multirow}                 % Multi-row table cells
\usepackage{caption}                  % caption formatting
\usepackage{subcaption}               % Sub-figures and sub-tables

% Math enhancements (optional but recommended)
\usepackage{bm}                       % Bold math symbols

% Reference enhancements (optional but recommended)
\usepackage{cleveref}                 % Intelligent cross-referencing
% Render DOIs as inline hyperlinks (not footnotes).  The doi.sty package
% is intentionally omitted because it renders DOIs as footnotes.  This
% override typesets them as clickable links within the bibliography entry.
\providecommand{\doi}[1]{\href{https://doi.org/#1}{\nolinkurl{doi:#1}}}

% Configure figure numbering and captions
\renewcommand{\figurename}{Figure}
\captionsetup{
    justification=centering,
    font=small,
    labelfont=bf,
    labelsep=period
}

% Configure table numbering and captions
\renewcommand{\tablename}{Table}
\captionsetup[table]{
    justification=centering,
    font=small,
    labelfont=bf,
    labelsep=period
}

% Configure section numbering
\setcounter{secnumdepth}{3}
\renewcommand{\thesection}{\arabic{section}}
\renewcommand{\thesubsection}{\arabic{section}.\arabic{subsection}}
\renewcommand{\thesubsubsection}{\arabic{section}.\arabic{subsection}.\arabic{subsubsection}}

% Configure equation numbering
\numberwithin{equation}{section}

% Configure hyperref for proper linking
\hypersetup{
    colorlinks=true,
    linkcolor=red,
    citecolor=red,
    urlcolor=red,
    filecolor=red,
    anchorcolor=red,
    pdfborder={0 0 0},
    bookmarks=true,
    bookmarksnumbered=true,
    bookmarkstype=toc,
    pdftitle={Ento-Linguistics: Language, Ambiguity, and Scientific Communication in Entomology},
    pdfsubject={Scientific Terminology Analysis in Entomology},
    pdfproducer={XeLaTeX}
    % pdfauthor and pdfkeywords are set dynamically in _title_page.tex
    % from config.yaml to avoid metadata drift
}

% ============================================================================
% PACKAGE CONFIGURATION
% ============================================================================

% Configure cleveref for intelligent cross-references
\crefname{section}{Section}{Sections}
\crefname{subsection}{Subsection}{Subsections}
\crefname{subsubsection}{Subsubsection}{Subsubsections}
\crefname{equation}{Equation}{Equations}
\crefname{figure}{Figure}{Figures}
\crefname{table}{Table}{Tables}
\crefname{appendix}{Appendix}{Appendices}

% Configure fonts for Unicode support with fallbacks
\usepackage{newunicodechar}
%\newunicodechar definitions for Unicode characters used in the manuscript
\newunicodechar{⁴}{\textsuperscript{4}}
\newunicodechar{₄}{\textsubscript{4}}
\newunicodechar{²}{\textsuperscript{2}}
\newunicodechar{₀}{\textsubscript{0}}
\newunicodechar{₁}{\textsubscript{1}}
\newunicodechar{₂}{\textsubscript{2}}
\newunicodechar{₃}{\textsubscript{3}}
% Arrow characters used in abstract, conclusion, and supplemental methods
\newunicodechar{→}{\ensuremath{\rightarrow}}
\newunicodechar{←}{\ensuremath{\leftarrow}}
\newunicodechar{↔}{\ensuremath{\leftrightarrow}}
% Additional symbols
\newunicodechar{≥}{\ensuremath{\geq}}
\newunicodechar{≤}{\ensuremath{\leq}}
\newunicodechar{≠}{\ensuremath{\neq}}
\newunicodechar{…}{\ldots}

% ============================================================================
% FONTS AND TYPOGRAPHY
% ============================================================================

% Use standard fonts for better compatibility
% NOTE: T1 fontenc is intentionally omitted — it is incompatible with xelatex
% (which uses Unicode natively). Including it generates font warnings that
% can cascade into compilation failures.
\usepackage{lmodern}

% ============================================================================
% CODE BLOCK STYLING
% ============================================================================

% code block styling for better contrast and readability
\usepackage{fancyvrb}
\usepackage{fvextra}   % Extends fancyvrb: adds bgcolor, breaklines, etc.
\usepackage{xcolor}
\usepackage{listings}

% Define custom colors for code blocks
\definecolor{codebg}{RGB}{248, 248, 248}      % Very light gray background
\definecolor{codeborder}{RGB}{200, 200, 200}  % Medium gray border
\definecolor{codefg}{RGB}{34, 34, 34}         % Dark gray text
\definecolor{commentcolor}{RGB}{102, 102, 102} % Comment color
\definecolor{keywordcolor}{RGB}{0, 0, 0}       % Keyword color
\definecolor{stringcolor}{RGB}{0, 102, 0}      % String color

% Configure Verbatim environment for inline code
\DefineVerbatimEnvironment{Verbatim}{Verbatim}{%
    fontsize=\small,
    breaklines=true,
    breakanywhere=true,
    frame=single,
    framerule=0.5pt,
    framesep=3pt,
    rulecolor=\color{codeborder},
    bgcolor=codebg
}

% Configure code block styling
\DefineVerbatimEnvironment{Highlighting}{Verbatim}{%
    fontsize=\footnotesize,
    breaklines=true,
    breakanywhere=true,
    frame=single,
    framerule=0.5pt,
    framesep=5pt,
    rulecolor=\color{codeborder},
    bgcolor=codebg
}

% Style inline code with \texttt
\DeclareRobustCommand{\codeunderscore}{\textunderscore\allowbreak}
\renewcommand{\texttt}[1]{%
    {\color{codefg}\ttfamily\let\_\codeunderscore #1}%
}

% Configure listings package for code blocks
\lstset{
    backgroundcolor=\color{codebg},
    basicstyle=\footnotesize\ttfamily\color{codefg},
    breakatwhitespace=false,
    breaklines=true,
    captionpos=b,
    commentstyle=\color{commentcolor},
    deletekeywords={...},
    escapeinside={\%*}{*)},
    extendedchars=true,
    frame=single,
    framerule=0.5pt,
    framesep=5pt,
    keepspaces=true,
    keywordstyle=\color{keywordcolor}\bfseries,
    language=Python,
    morekeywords={*,...},
    numbers=left,
    numbersep=5pt,
    numberstyle=\tiny\color{codefg},
    rulecolor=\color{codeborder},
    showspaces=false,
    showstringspaces=false,
    showtabs=false,
    stepnumber=1,
    stringstyle=\color{stringcolor},
    tabsize=4,
    title=\lstname
}

% Override any Pandoc default lstset configurations
\AtBeginDocument{
    \lstset{
        backgroundcolor=\color{codebg},
        basicstyle=\footnotesize\ttfamily\color{codefg},
        frame=single,
        framerule=0.5pt,
        framesep=5pt,
        rulecolor=\color{codeborder},
        numbers=left,
        numbersep=5pt,
        numberstyle=\tiny\color{codefg}
    }
}

% Configure bibliography
% Note: Using plainnat with natbib package for proper citation processing
% The bibliography style and commands (\bibliographystyle and \bibliography) are in 99_references.md
\setlength{\bibsep}{2pt}

% Simple page break support for document structure
% Note: Page breaks are handled in the markdown generation, not here

% ============================================================================
% DOCUMENT FORMATTING
% ============================================================================

% Ensure proper spacing and formatting
\frenchspacing  % Single space after periods
\linespread{1.2}  % Slightly increased line spacing for readability

% ============================================================================
% TITLE PAGE OVERRIDE
% ============================================================================
% The title page, author block, DOI, and PDF metadata are generated at
% build time by _render_pdf_override.py from config.yaml, written to
% _title_page.tex, and \input'd here.  This avoids hardcoding metadata
% in two places.  The generated file redefines \maketitle using \parbox
% blocks instead of the default tabular (which breaks under TeXLive 2026).
\IfFileExists{_title_page.tex}{% Auto-generated from config.yaml — do not edit by hand.
% Uses a custom \maketitle to avoid TeXLive 2026 tabular/\and breakage.
\makeatletter
\renewcommand{\maketitle}{%
  \begin{center}
    {\LARGE \@title \par}
    \vskip 1.5em
    \parbox[t]{0.45\textwidth}{\centering
      \textbf{Daniel Ari Friedman}\\
      {\small Active Inference Institute, APOIDEAS}\\
      {\footnotesize {\ttfamily daniel@activeinference.institute}\\ ORCID:~0000-0001-6232-9096}
    }
    \hfill
    \parbox[t]{0.45\textwidth}{\centering
      \textbf{Tucker Cahill Chambers}\\
      {\small APOIDEAS}\\
      {\footnotesize ORCID:~0009-0008-3793-7872}
    }
    \vskip 1em
    2026-10-08 – Version 1.3.0
  \end{center}
  \par\vskip 1.5em
}
\makeatother

% PDF metadata (synced from config.yaml keywords and authors)
\hypersetup{%
  pdfauthor={Daniel Ari Friedman and Tucker Cahill Chambers},
  pdfkeywords={entomology, scientific terminology, language analysis, discourse analysis, ant biology, eusociality, conceptual frameworks, active inference, terminology networks, Markov blankets, stigmergy, complex systems, complexity, representation sensitivity, computational linguistics},
}
}{}
\usepackage{bookmark}
\IfFileExists{xurl.sty}{\usepackage{xurl}}{} % add URL line breaks if available
\urlstyle{same}
% fallback for those not using the hyperref driver hyperxmp:
\makeatletter
\@ifundefined{xmpquote}{\newcommand{\xmpquote}[1]{#1}}{}
\makeatother
\hypersetup{
  pdftitle={Ento-Linguistics: Language{,} Ambiguity{,} and Scientific Communication in Entomology},
  pdfauthor={Daniel Ari Friedman; Tucker Cahill Chambers},
  colorlinks=true,linkcolor=red,urlcolor=red,citecolor=red,anchorcolor=red,filecolor=red,
  pdfcreator={LaTeX via pandoc}}

\title{Ento-Linguistics: Language, Ambiguity, and Scientific
Communication in Entomology}
\usepackage{etoolbox}
\makeatletter
\providecommand{\subtitle}[1]{% add subtitle to \maketitle
  \apptocmd{\@title}{\par {\large #1 \par}}{}{}
}
\makeatother
\subtitle{Complex Systems, Measurement, and a Six-Domain Framework}
\author{Daniel Ari Friedman \and Tucker Cahill Chambers}
\date{2026-10-08 -- Version 1.3.0}

\begin{document}
\maketitle

{
\setcounter{tocdepth}{2}
\tableofcontents
}
\section{Abstract}\label{sec:abstract}

Terms such as \emph{queen}, \emph{worker}, and \emph{colony} connect
biological descriptions to familiar social concepts. This study
introduces a six-domain Ento-Linguistic framework and an open-source
descriptive text-analysis pipeline. The headline layer contains 7599
source-identified abstracts from 7609 stored strings; 10 unreconciled
strings are retained but excluded. This headline layer contains 998,645
processed tokens and 11,700 candidate terms, of which 1330 receive
rule-based domain assignments. The pipeline separates observed
document-level term co-occurrence from a map of 6 predefined concept
categories with 15 vocabulary-overlap relationships. Among assigned
terms, 11.9\% receive multiple labels; this measures classification
overlap rather than semantic drift. Complementary analyses use 6997 PMC
full-text records, 2430 historical OCR documents, and a separate arXiv
layer. TF-IDF clustering and Shannon entropy summarize sentence-context
distributions, while lexical patterns identify candidate framing
contexts. A separately receipted fixed-margin comparison finds fewer
edges and lower clustering in the selected-term graph than in its
finite-chain randomized reference, indicating concentration of
co-occurrence relative to the preserved incidence margins. Clarity,
Appropriateness, Consistency, and Evolvability (CACE) are proposed as
heuristic evaluation dimensions. Neither cluster entropy nor marker
occurrence establishes distortion of biological understanding, and CACE
scores have not been validated against independent human judgments. A
deterministic vocabulary--threshold diagnostic makes the dependence of
density and clustering on network representation explicit. Complexity
motivates questions about units, coupling, feedback and scale; it is not
inferred from a dense lexical projection. Provenance gaps, mixed-topic
retrieval, OCR errors, overlapping groups, and explicitly bounded
analyses limit interpretation. The contribution is a reproducible
descriptive workflow and a framework for subsequent annotated,
hypothesis-driven research, rather than a causal test of language
shaping scientific thought. Code and data lineage:
https://github.com/docxology/ento\_linguistics.

\begin{figure}[htbp]
\centering
\includegraphics[width=\textwidth,height=0.65\textheight,keepaspectratio]{/Volumes/GitVault/Local/worktrees/ento-v1.3.1/ento_linguistics/output/extensions/network_reading/graphical_abstract.png}
\caption{Graphical abstract. Biological processes, scientific discourse and computational representations have distinct observational units. The six overlapping domains organize questions; observed textual proxies and proposed validation designs retain separate evidential roles. Arrows in the biological lane are schematic, not a fitted model or newly measured interaction network.}
\label{fig:graphical_abstract}
\end{figure}

\clearpage

\section{Introduction}\label{sec:introduction}

\subsection{Terminology as a Research
Question}\label{terminology-as-a-research-question}

Scientific terminology connects observations, categories, and
explanatory models. In social-insect research, familiar labels such as
queen, worker, caste, kin, and colony can serve technical functions
while retaining ordinary-language connotations. The relevant question is
when those connotations help communication, when explicit definitions
are needed, and how either possibility can be studied without inferring
beliefs from word counts alone.

Philosophical and linguistic work supplies motivation for examining
terminology within scientific practice
\citep{latour1987, longino1990, lakoff1980metaphors}. Entomological
discussions of categories and loaded metaphors provide field-specific
context \citep{gordon1992wittgenstein, herbers2006, herbers2007}. These
perspectives motivate empirical questions; they do not make every
conventional term misleading or demonstrate that a replacement improves
scientific understanding.

\subsection{The Challenge of Terminological
Reform}\label{the-challenge-of-terminological-reform}

Established terms support literature discovery and can carry precise
operational meanings. Proposed alternatives therefore require
context-specific comparison rather than automatic replacement. A
task-based description may be useful for a behavioral observation
without replacing a developmentally defined caste category. Molecular
and developmental research makes that distinction especially important
\citep{sumner2018molecular, qiu2022canalized}.

CACE---Clarity, Appropriateness, Consistency, and Evolvability---is
introduced as an explicit set of evaluation questions. Its numerical
scoring rules are one inspectable implementation, rather than an
independently validated measure of understanding. The biological
mechanism, intended referent, reader population, and definition supplied
by an author remain central to evaluating terminology.

\subsection{From Labels to Biological
Explanation}\label{from-labels-to-biological-explanation}

A useful terminology analysis separates the expression, its referent,
and the mechanism invoked to explain that referent. \emph{Worker} can
identify a developmental caste (typically non-reproductive females) or
describe an individual performing a task; neither usage alone specifies
how task allocation is regulated. Definitions of caste have themselves
been examined as a conceptual problem \citep{villet1992caste}, while
experimental work shows that individual experience can contribute to
persistent division of labor \citep{ravary2007}. These examples locate
the research question in the relationship between definitions and
evidence, rather than in a presumption that social vocabulary is
necessarily inaccurate.

The six domains provide questions to ask of a passage: what entity is
described, at which biological scale, using which observable criteria,
and with what explanatory commitment? Frequency and co-occurrence help
locate passages for this inquiry. Assessing whether a label obscures a
mechanism requires examining those passages and evaluating readers'
inferences. This distinction makes the framework useful for comparative
reading without treating the taxonomy as a discovered ontology.

\subsection{Six Analytical Domains}\label{six-analytical-domains}

The framework defines six overlapping domains:

\begin{enumerate}
\def\labelenumi{\arabic{enumi}.}
\tightlist
\item
  \textbf{Unit of Individuality:} terms identifying ants, nestmates,
  colonies, collectives, and superorganisms; how an author specifies the
  unit being observed or modeled.
\item
  \textbf{Behavior and Identity:} distinctions between a task, an
  observed behavior, a persistent propensity, and a developmental
  category.
\item
  \textbf{Power and Labor:} queen, worker, caste, and related labels;
  whether descriptions imply control mechanisms or simply name
  biological roles.
\item
  \textbf{Sex and Reproduction:} reproductive and developmental
  categories; how definitions accommodate variation across taxa and
  mating systems.
\item
  \textbf{Kin and Relatedness:} pedigrees, demographic assumptions, and
  relatedness terminology. For example, expected sister relatedness of
  \(r=0.75\) requires an outbred haplodiploid pedigree with one mother
  and one haploid father; it is not a universal property of colonies.
\item
  \textbf{Economics:} resource, allocation, investment, and related
  expressions; distinctions between measured energetic expenditure,
  fitness consequences, and metaphorical usage.
\end{enumerate}

These are predefined organizational categories. A term receiving several
labels does not establish that its meaning changed over time. Each
domain also requires qualitative annotation before lexical patterns can
be interpreted as ambiguity or framing.

\subsection{Research Approach}\label{research-approach}

The headline computation analyzes 7599 source-identified abstracts,
yielding 998,645 processed tokens and 11,700 candidates, of which 1330
receive domain assignments. Complementary PMC, BHL, and arXiv layers
retain separate source identities. The workflow distinguishes corpus
frequencies, observed document co-occurrence, vocabulary overlap,
context-cluster entropy, and heuristic scores.

Active Inference and multiscale modeling provide a theoretical
perspective \citep{friston2010free, friedman2021active}, but the present
pipeline does not fit a generative model of scientific language or
measure variational free energy. A Markov blanket specifies
conditional-independence relationships in a model; it is not a lexical
security filter or a biological boundary established by terminology
alone.

The contribution is a descriptive workflow, a proposed taxonomy, and
explicit questions for subsequent validation. Historical OCR frequencies
provide dated source-layer observations rather than proof of conceptual
origins or causal reform. Methods and limitations state source custody
gaps, convenience-sample retrieval, statistical dependence, and bounded
analyses so that interpretation remains tied to the evidence actually
produced.

Complex systems supplies an organizing question: which units, relations,
feedbacks and timescales does a scientific passage specify? Network
representations and dynamics address different problems
\citep{newman2003structure}, while different complexity measures require
different targets and assumptions \citep{ladyman2013complex}. This study
therefore keeps observed biological mechanisms, lexical descriptions and
computational proxies distinct. The v1.3.0 representation diagnostic
makes vocabulary and edge-threshold dependence inspectable; separate
validation is required before transferring those measures to meaning,
communication or biological organization.

\clearpage

\section{Methods}\label{sec:methodology}

The workflow separates acquisition, descriptive analysis, and
theoretical interpretation. Supplemental Methods
\ref{sec:supplemental_methods} and \ref{sec:supplemental_infrastructure}
describe the implementation. Section \ref{sec:supplemental_analysis}
develops conceptual proposals rather than fitted generative models.

\subsection{Data Acquisition}\label{sec:data_acquisition}

\subsubsection{Search Strategy and Source
Selection}\label{search-strategy-and-source-selection}

The archived input is the ordered abstract list in
\emph{data/corpus/abstracts.json}. The headline analysis selects strings
whose digest maps to an identified PubMed record. Unreconciled strings
remain archived but are excluded, without altering their bytes. The
acquisition script delegates to \emph{src/pipeline/corpus\_build.py};
its base queries and the named growth queries in
\emph{src/data/literature\_mining.py} are the executable query
definitions. Digest-indexed provenance records PMIDs, DOIs, titles,
years, journals, and query names where available. This corpus
accumulated across retrieval runs; a single illustrative query does not
reconstruct its history.

Three complementary layers remain separate: PMC title, abstract, and
body text in \emph{data/fulltexts/} shards; BHL mirror OCR in
\emph{data/bhl/} shards; and arXiv title/abstract records in
\emph{data/corpus/arxiv\_records.json}. The arXiv harvester uses
category-restricted queries over q-bio.PE and nlin.AO. Neither full
texts nor preprints are merged into the headline abstract list. Query
and retrieval metadata are retained in source sidecars.

\subsubsection{Corpus Composition and
Cleaning}\label{corpus-composition-and-cleaning}

{\def\LTcaptype{none} % do not increment counter
\begin{longtable}[]{@{}ll@{}}
\toprule\noalign{}
Headline metric & Value \\
\midrule\noalign{}
\endhead
\bottomrule\noalign{}
\endlastfoot
Stored abstract strings & 7609 \\
Source-identified analyzed abstracts & 7599 \\
Unreconciled strings excluded & 10 \\
Processed tokens & 998,645 \\
Unique token types & 47,216 \\
Candidate terms & 11,700 \\
Domain-assigned terms & 1330 \\
\end{longtable}
}

These are archived-input and selected-corpus counts, not estimates of
all entomological literature. Broad searches retrieve adjacent
biological topics and algorithmic uses of ant terminology; a keyword
match does not establish relevance to ant biology. Historical volumes
can contain substantial non-entomological material. Of 7073 stored PMC
records, 76 whose titles mark them as correction, erratum or retraction
notices or as retracted items remain archived but are excluded from
analysis. This removes the standardized repeated notice text. A title
rule is not a complete article-type or retraction-status screen, and
relevance screening is not complete. The custody audit in
\emph{output/reports/corpus\_audit.json} records missing digest
provenance, repeated text, identifier mismatches, and unused sidecar
entries. Gaps are retained and disclosed; undocumented source identities
are not inferred. Exact text-digest matches to retrieved PubMed
abstracts can recover metadata without rewriting text. This establishes
source identity rather than ant-topic relevance. Digest-keyed sidecars
retain only one metadata entry when several records contain identical
text. Relevance annotation and complete source reconciliation are needed
before treating a layer as a curated field-wide sample.

TextProcessor normalizes text, tokenizes with NLTK, filters punctuation
and stop words, and optionally lemmatizes with WordNet. Extraction
disables lemmatization when counting surface-form terms. NLTK datasets
are installed separately from Python dependencies; selected English
tokenizer, stopword and WordNet contents are bound into analysis
signatures and receipts. Article-level licenses require individual
checking; indexing in PubMed/PMC does not establish reuse permission.

\subsubsection{Domain Coverage
Verification}\label{domain-coverage-verification}

The six vocabularies are defined by TerminologyExtractor.DOMAIN\_SEEDS:
Unit of Individuality, Behavior and Identity, Power and Labor, Sex and
Reproduction, Kin and Relatedness, and Economics. Coverage means that
extraction assigned terms to these labels; it does not measure retrieval
recall or validate the taxonomy against independent annotation.

\subsection{Statistical Analysis}\label{sec:statistical_analysis}

\subsubsection{Term Extraction and
Classification}\label{term-extraction-and-classification}

Extraction counts candidate tokens, retains configured seed words
subject to length/frequency and preprocessing rules, applies remaining
lexical-pattern filters, and assigns domains through direct seed
matching, compound-word overlap, and fallback patterns. Three-token
windows provide up to thirty deduplicated short contexts for heuristic
scoring. They do not expand assignments by co-occurrence. Candidate
filtering also admits scientific substrings and compound tokens, so
candidates can include unrelated words and OCR artifacts. Candidate
inclusion does not itself assign a domain or establish biological
relevance. An n-gram utility is available but the headline extractor
does not invoke it to extract multi-word phrases. The BHL
literal-frequency pass separately matches multi-word seed phrases as
consecutive tokens.

Minimum frequencies differ: one occurrence for headline abstracts,
twenty for PMC, and two for arXiv and the BHL stack. These choices
affect vocabulary size and prevent interpreting raw term counts as
directly comparable measures across layers.

\subsubsection{Semantic Entropy}\label{semantic-entropy}

Terms with at least five usable sentence contexts undergo TF-IDF
vectorization and seeded KMeans clustering (random state 42, ten
initializations). Shannon entropy measures occupancy:

\begin{equation}\label{eq:semantic_entropy}
H(t)=-\sum_{i=1}^{k}p_i\log_2p_i.
\end{equation}

The requested cluster count is
\(k=\max(2,\min(5,n-1,\max(2,\lfloor\sqrt{n}\rfloor)))\); occupied
clusters may be fewer. Normalized entropy divides by
\(\log_2 k_{\mathrm{occupied}}\), with zero for a single occupied
cluster. Only successfully computed estimates enter domain means;
insufficient-context and computation failures are recorded as
exclusions. Domain figures and statistical descriptives share this
sentence-context calculation. The \(H>2\) bits flag is exploratory,
without independent calibration. Clusters are computational partitions,
not annotated word senses; high entropy does not by itself demonstrate
ambiguity or temporal semantic drift.

\subsubsection{Measurement Targets and
Validation}\label{sec:measurement_validation}

Each computational quantity has a defined observational unit.
Frequencies count processed occurrences; network edges count shared
documents; domain overlap counts classifier memberships; entropy
summarizes sentence-context cluster occupancy. Framing markers and CACE
scores apply configured rules to selected contexts and terms. None
directly measures an author's intention or a reader's understanding.
Problem-specific validation is a central requirement of automated text
analysis \citep{grimmer2013text}; reproducible calculation and valid
interpretation address different questions.

A future evaluation should freeze annotation rules before inspecting
score rankings and retain a held-out set of source contexts. Annotation
should distinguish biological referent, task versus developmental
category, technical definition, figurative use, and uncertainty.
Sampling should include frequent and rare terms, assigned and unassigned
candidates, and contexts without marker matches. These strata allow
false positives and false negatives to be examined, rather than
assessing only examples selected by the pipeline. Agreement,
adjudication rules, and performance by source layer should be reported.
This is a proposed validation design, not an executed component of the
current analysis.

\subsubsection{Domain-Level Statistical
Tests}\label{domain-level-statistical-tests}

Welch tests, Benjamini--Hochberg adjustments, small-sample corrected
standardized differences, and one-way ANOVA summarize valid per-term
entropies. Groups with fewer than two usable values are omitted. These
are exploratory comparisons. Terms can belong to multiple groups and
derive contexts from shared documents, violating the simple independence
interpretation. Multiplicity correction does not repair dependence,
retrieval bias, or unequal context counts. Threshold crossings do not
establish calibrated population significance or causal effects.
Confirmatory inference requires independent annotation and a
document-level sampling model accounting for shared terms and
overlapping domains.

\subsubsection{Conceptual Network
Analysis}\label{conceptual-network-analysis}

The conceptual map assigns terms to six predefined categories. Its 15
links summarize vocabulary overlap:

\begin{equation}\label{eq:overlap_coefficient}
w_{AB}=\frac{|A\cap B|}{\min(|A|,|B|)}.
\end{equation}

The terminology graph counts actual documents containing each pair among
the hundred most frequent domain-assigned terms, breaking frequency ties
lexically. Whole-word matching is case-insensitive; repeated mentions
within a document count once. Shared labels never create an observed
edge. Document co-occurrence does not establish a semantic or causal
relationship. The separate overlap heatmap summarizes classifier
assignments. A fixed-margin Curveball randomization of the binary
document--term incidence matrix supplies a conditional reference for the
projected edge count and mean local clustering (Supplemental Section
\ref{sec:fixed_margin}).

\subsubsection{Rhetorical and Discourse
Analysis}\label{rhetorical-and-discourse-analysis}

Regex and keyword analyzers identify candidate rhetorical,
argumentative, and framing expressions. Framing is evaluated in
three-token windows around extracted terms. Domain framing proportions
measure occurrence contexts containing an anthropomorphic marker;
multi-domain occurrences contribute to each domain but once to the
overall tally. These are lexical proxies without gold-standard
annotation validation. Discourse counts are accompanied by
analyzed-document counts and sampling fractions. PMC discourse uses a
deterministic approximately one-fifth sample of eligible texts rather
than a full-text census.

\subsubsection{CACE Evaluation}\label{cace-evaluation}

Four heuristic dimensions define CACE:

\begin{equation}\label{eq:cace_clarity}
\mathrm{Clarity}(t)=\min(1,\max(0,1-H(t)/3.32)),
\end{equation}

\begin{equation}\label{eq:cace_appropriateness}
\mathrm{Appropriateness}(t)=\max(0,1-[0.4\,\mathbf{1}_{t\in\mathcal A}+0.1|\mathrm{overlap}(t,\mathcal A)|+0.05\max(|D(t)|-1,0)]),
\end{equation}

\begin{equation}\label{eq:cace_consistency}
\mathrm{Consistency}(t)=\frac{2}{n(n-1)}\sum_{i<j}\frac{\mathbf x_i\cdot\mathbf x_j}{\|\mathbf x_i\|\|\mathbf x_j\|},
\end{equation}

\begin{equation}\label{eq:cace_evolvability}
\mathrm{Evolvability}(t)=\tfrac12\min(1,|D(t)|/3)+\tfrac12\min(1,|S_t|/3).
\end{equation}

The implemented Clarity normalization constant is 3.32, approximately
\(\log_2 10\); it is a design choice rather than the five-cluster
ceiling. Here \(\mathcal A\) is the configured anthropomorphic
vocabulary, \(\mathbf x_i\) the TF-IDF vector of the \(i\)-th of \(n\)
short extraction contexts, \(D(t)\) the domain assignments, and \(S_t\)
scale-marker categories found in contexts. The aggregate is their
arithmetic mean. Domain CACE aggregates and their figure use the same
up-to-fifty-term sample, ordered by decreasing frequency and lexical
tie-break, and at most ten short extraction contexts per term. Computed
entropy enters Clarity where available. Default zero entropy for
unavailable estimates and Consistency of 0.5 for fewer than two contexts
are conventions, not evidence of clarity or consistency. Weights and
thresholds are design choices; no completed coefficient sensitivity
study, inter-rater reliability measurement, or independent predictive
validation is claimed.

\subsubsection{Historical Analysis and
Coverage}\label{historical-analysis-and-coverage}

All 2430 stored BHL documents enter literal seed-term frequencies per
10,000 OCR tokens by era. Literal spellings, hyphens, and OCR errors
affect rates. The computational stack streams extraction and framing
over all stored documents by default. Coverage metadata record available
and analyzed documents and characters. Entropy is restricted to twenty
frequent candidates per era, drawing contexts from all era documents. An
optional positive development character budget selects complete
documents by deterministic dyadic traversal across shard order; its
fingerprint and coverage metadata distinguish it from a complete run.
Neither the full stored corpus nor a bounded development selection is
presumed representative of historical entomology.

\subsubsection{Validation and
Reproducibility}\label{validation-and-reproducibility}

Tests exercise numerical examples, real files, corpus slices, and local
HTTP servers. Negative controls reject malformed text, stale content,
missing receipts, and failed renderer commands. They establish software
behavior, not scientific validity of the proxy measures. Python
dependencies are resolved by uv.lock; NLTK resources are separate
prerequisites. Content fingerprints bind analyses to ordered records,
project source, and locked dependencies. A completed manifest binds
input and output files; rendering rejects changed artifacts and nonzero
toolchain exits.

\subsection{Network Representation
Sensitivity}\label{sec:network_reading}

Using the verified document--term incidence and stable frequency
ranking, we examine nested vocabularies of 25, 50, 75 and 100 terms at
inclusive shared-document thresholds of 1, 2, 5, 10 and 20. This
declared representation grid complements the fixed-margin reference; it
is neither a tuned model nor a random sample. Projection and dynamics
address distinct questions \citep{newman2003structure}.

Every selected term remains a node, including isolates. An edge exists
when \(w(u,v)\geq\tau\), and density is
\(\rho_{n,\tau}=2|E_{n,\tau}|/[n(n-1)]\). Mean local unweighted
clustering is zero for degree-zero/one nodes. Component count,
largest-component size, isolates and documents containing any selected
term are exported. Changing vocabulary size changes both possible pairs
and selected incidences; density across sizes need not be monotone.

The separate receipt binds inputs, implementation, lock and all outputs.
Validation reconstructs the baseline and replays the report. No
additional null model, bootstrap, uncertainty interval or community
optimization is executed. The statistics diagnose representation
dependence, not semantic validity or biological complexity.

\clearpage

\section{Results: Corpus Analysis and Terminology
Networks}\label{sec:experimental_results}

\subsection{Terminology Extraction Across
Domains}\label{terminology-extraction-across-domains}

Analysis of the source-identified headline layer yields 11,700 candidate
terms from 7599 abstracts and 998,645 processed tokens. Of these
candidates, 1330 receive domain assignments. These counts exclude 10
unreconciled archived strings; broad retrieval still limits relevance
and representativeness.

\begin{table}[h]
\centering
\begin{tabular}{|l|c|c|c|}
\hline
\textbf{Domain} & \textbf{Terms} & \textbf{Frequency} & \textbf{Bridging terms} \\
\hline
Unit of Individuality & 530 & 28,194 & 24 \\
Behavior \& Identity & 262 & 17,119 & 91 \\
Power \& Labor & 295 & 12,533 & 139 \\
Sex \& Reproduction & 245 & 9334 & 62 \\
Kin \& Relatedness & 84 & 4040 & 4 \\
Economics & 79 & 4508 & 3 \\
\hline
\end{tabular}
\caption{Rule-based domain assignments and corpus frequencies. A term can receive several labels, so domain counts and frequencies are not mutually exclusive. Bridging means multiple labels, not an observed transfer of meaning.}
\label{tab:terminology_extraction}
\end{table}

Among assigned terms, 11.9\% have multiple domain labels. This is a
property of the classifier and corpus. Temporal semantic drift would
require explicit time-indexed meaning comparison, while lexical,
contextual, and scale ambiguity require independent sense annotation or
validated proxies. Those measurements are not established by label
overlap.

The processed vocabulary has type-token ratio 0.0473. Its most frequent
WordNet lemmas are ant (20,494), specie (10,620), and colony (7436).
Lemmatization reduces inflected forms and can truncate words such as
\emph{species} to \emph{specie}; extracted terms retain surface forms
(Supplemental Section \ref{sec:supplemental_methods}). Frequency
identifies recurring lexical material; it does not establish its
conceptual importance or the intentions of authors.

\subsection{Terminology Network
Structure}\label{terminology-network-structure}

The observed terminology graph uses the hundred most frequent
domain-assigned terms. Its edge weight counts documents containing both
terms:

\begin{equation}\label{eq:network_edge_weight}
w(u,v)=\sum_{d=1}^{N}\mathbf{1}[u\in d]\mathbf{1}[v\in d].
\end{equation}

Whole-word matches are case-insensitive, and repeated mentions within a
document do not add weight. No edge is inferred from shared labels or
extraction order. The graph has mean local (unweighted) clustering
coefficient 0.8955; this statistic does not measure conceptual
coherence, communication quality, or resistance to reform.

\begin{figure}[h]
\centering
\includegraphics[width=0.95\textwidth,height=0.78\textheight,keepaspectratio]{/Volumes/GitVault/Local/worktrees/ento-v1.3.1/ento_linguistics/output/figures/terminology_network.png}
\caption{Observed document co-occurrence among the hundred highest-frequency domain-assigned terms in the 7599-abstract layer. Nodes represent terms, node area uses a square-root frequency scale, color identifies the primary domain, and edge width scales shared-document counts to a bounded display range (Eq.~\ref{eq:network_edge_weight}). Isolated nodes are omitted from the display and up to twenty frequent terms are considered for collision-filtered labels. The right panel ranks the ten strongest observed pairs by shared-document count, with lexical tie-breaking. All observed edges remain in the left graph, whose unweighted layout reduces visual concentration around frequent pairs. Layout and dense regions have no causal or hierarchical interpretation.}
\label{fig:terminology_network}
\end{figure}

\subsubsection{Comparison with a Fixed-Margin
Reference}\label{comparison-with-a-fixed-margin-reference}

Density and clustering depend on how many selected terms each abstract
contains and how many abstracts contain each term. A Curveball
randomization of the binary document--term incidence matrix preserves
both margins (Supplemental Section \ref{sec:fixed_margin}). Over 3
chains and 600 retained draws for the 7599 source-identified abstracts
and 100 terms, the randomized reference averages 4533.6 edges (central
95\% of draws 4501--4568) and mean local clustering 0.934
(0.930--0.939). The observed graph has 4235 edges and clustering 0.895,
below both envelopes (Figure \ref{fig:network_robustness}). A
sensitivity protocol with longer burn-in, wider spacing and different
seeds (300 draws) gives reference means of 4534.5 edges and 0.934
clustering. Between-chain \(\hat R\) is 1.006 for edges and 1.003 for
clustering; this diagnostic does not prove mixing.

Selected terms therefore co-occur in fewer distinct pairs, with less
closed triadic structure, than their margins alone would produce:
co-occurrence is concentrated among particular pairs. The envelopes
describe the sampled conditional reference rather than uncertainty in
the observed values, and the comparison does not identify a linguistic
cause.

\begin{figure}[htbp]
\centering
\includegraphics[width=\textwidth,height=0.7\textheight,keepaspectratio]{/Volumes/GitVault/Local/worktrees/ento-v1.3.1/ento_linguistics/output/extensions/network_robustness/network_robustness.png}
\caption{Fixed-margin comparison for the abstract terminology network. Top: retained-chain traces of projected edge count and mean local clustering for 3 seeds. Bottom: distributions of the 600 retained draws, with the observed value marked. Every draw preserves each abstract's selected-term count and each term's document frequency. The reference is conditional on the frequency-ranked vocabulary and the convenience corpus.}
\label{fig:network_robustness}
\end{figure}

\subsubsection{Vocabulary and Edge-Threshold
Diagnostic}\label{vocabulary-and-edge-threshold-diagnostic}

Across the declared 20 vocabulary--threshold combinations on the same
7599 identified abstracts, the 100-term graph at a one-document
threshold has density 0.8556. Requiring at least twenty shared documents
retains 1211 edges and gives density 0.2446. Figure
\ref{fig:network_reading} displays the full grid, including mean local
clustering with isolates retained. These are deterministic descriptions
of selected representations, not uncertainty intervals or evidence that
one threshold recovers true semantic relations. Clustering need not
decrease when edges are removed: changing neighborhoods also changes
local triangle denominators. The curves separate these topological
summaries rather than interpreting either as biological complexity.

\begin{figure}[htbp]
\centering
\includegraphics[width=\textwidth,height=0.65\textheight,keepaspectratio]{/Volumes/GitVault/Local/worktrees/ento-v1.3.1/ento_linguistics/output/extensions/network_reading/network_reading.png}
\caption{Vocabulary and minimum-shared-document sensitivity. Curves use nested frequency-ranked vocabularies and inclusive thresholds. The denominator includes all selected pairs, and clustering includes isolates as zero. Each point is recomputed from the same identified abstract layer; no new randomization or biological-network inference is represented.}
\label{fig:network_reading}
\end{figure}

Domain-assignment overlap is a different quantity, displayed separately
in Figure \ref{fig:domain_overlap}.

\begin{figure}[h]
\centering
\includegraphics[width=0.9\textwidth,height=0.78\textheight,keepaspectratio]{/Volumes/GitVault/Local/worktrees/ento-v1.3.1/ento_linguistics/output/figures/domain_overlap_heatmap.png}
\caption{Szymkiewicz--Simpson overlap coefficients between domain-assigned vocabularies (Eq.~\ref{eq:overlap_coefficient}). Each cell shows shared terms as a percentage of the smaller vocabulary; the diagonal is 100 by definition. Values reflect the current lexical classifier; observed zeros do not establish conceptual isolation.}
\label{fig:domain_overlap}
\end{figure}

\begin{figure}[h]
\centering
\includegraphics[width=0.9\textwidth,height=0.78\textheight,keepaspectratio]{/Volumes/GitVault/Local/worktrees/ento-v1.3.1/ento_linguistics/output/figures/domain_comparison.png}
\caption{Six descriptive panels show distinct term counts, mean extraction confidence, total frequency, mean successfully computed sentence-context entropy, bridging counts, and heuristic CACE means over up to fifty selected terms per domain. Extraction confidence is a configured score rather than calibrated classification accuracy. Entropy and CACE sample definitions are specified in Methods; missing entropy does not establish zero ambiguity.}
\label{fig:domain_comparison}
\end{figure}

\begin{figure}[htbp]
\centering
\includegraphics[width=0.9\textwidth,height=0.78\textheight,keepaspectratio]{/Volumes/GitVault/Local/worktrees/ento-v1.3.1/ento_linguistics/output/figures/concept_map.png}
\caption{Six predefined concept categories connected by vocabulary overlap. Node size summarizes associated terms, which can appear in several categories; the subtitle totals term associations rather than unique terms. Edge weights are classifier-defined overlaps, not observed causal connections or a discovered biological ontology.}
\label{fig:concept_map}
\end{figure}

\subsection{Framing Analysis}\label{framing-analysis}

Lexical markers identify contexts for further qualitative examination.
They do not distinguish metaphor from technical usage, demonstrate
author bias, or establish a language-induced change in biological
models. The same distinction applies to the domain-specific
interpretations in Section \ref{sec:domain_findings}.

\begin{figure}[h]
\centering
\includegraphics[width=0.9\textwidth,height=0.78\textheight,keepaspectratio]{/Volumes/GitVault/Local/worktrees/ento-v1.3.1/ento_linguistics/output/figures/anthropomorphic_framing.png}
\caption{Observed framing-marked terminology by canonical domain. Top: distinct extracted terms with at least one occurrence context matching an anthropomorphic pattern. Bottom: up to six terms per domain, selected by decreasing matched-context proportion, context count, and lexical ordering. Counts are neither curated vocabulary sizes nor occurrence frequencies. Occurrence-context framing proportions are exported separately.}
\label{fig:anthropomorphic}
\end{figure}

\clearpage

\section{Results: Domain-Specific Findings}\label{sec:domain_findings}

The six domains organize descriptive outputs and questions for
qualitative interpretation. Their assignments follow the configured
vocabulary and lexical rules. Domain-specific statements below do not
infer cognitive bias, biological mechanisms, or author intent from term
counts.

\subsection{Unit of Individuality}\label{unit-of-individuality}

This vocabulary includes references to individuals, colonies, and
collective organization. Its mean computed sentence-context entropy is
1.55 bits. Frequency and word-formation panels describe extracted
surface forms, while the scale panel counts keyword matches in term
names. These counts do not estimate biological boundaries. Figures
\ref{fig:domain_overview_grid}, \ref{fig:domain_patterns_grid}, and
\ref{fig:unit_individuality_patterns} support inspection of the
extraction.

\subsection{Power \& Labor}\label{power-labor}

Power and Labor contains 295 assigned terms and 139 terms with multiple
labels. Its occurrence-context anthropomorphic-marker proportion is
1.2\%, rather than a percentage of authors or publications using
misleading language. Discussions of loaded terminology motivate
contextual examination \citep{herbers2006, herbers2007}. Molecular work
on caste \citep{sumner2018molecular} and developmental canalization
\citep{qiu2022canalized} also make it important to distinguish
developmental phenotypes from temporary task categories.

Figures \ref{fig:power_labor_frequencies} and
\ref{fig:power_labor_ambiguities} show observed term frequency and
context-cluster entropy. Figure \ref{fig:concept_hierarchy} ranks
predefined concept categories by summed vocabulary overlap; it does not
show a biological hierarchy or term-level betweenness.

\subsection{Behavior \& Identity}\label{behavior-identity}

The mean computed sentence-context entropy is 1.74 bits. Task labels
provide useful candidates for context annotation. Evidence of behavioral
flexibility in ants \citep{ravary2007, gordon2010} motivates asking when
a label denotes an observation, a persistent propensity, or a
morphological category. The corpus statistics alone do not establish
that categorical labels obscure that flexibility.

\subsection{Sex \& Reproduction}\label{sex-reproduction}

This domain groups reproductive and developmental terminology. Its mean
computed entropy is 1.73 bits. The presence of paired labels does not by
itself demonstrate a conceptual opposition. Reproductive systems and
caste development vary across taxa, and the epigenetic review
\citep{oldroyd2021epigenetics} provides background rather than
validation of a lexical classifier.

\subsection{Kin \& Relatedness}\label{kin-relatedness}

The mean computed entropy is 1.76 bits. Relatedness terms require
biological and demographic context: numerical coefficients depend on
pedigrees, mating systems, and population structure. The present outputs
do not estimate those quantities.

\subsection{Economics}\label{economics}

The classifier assigns 79 terms to Economics, with 3 receiving multiple
labels. Its mean computed entropy is 1.85 bits. Low overlap can follow
directly from lexicon design; a zero overlap does not establish a closed
conceptual subsystem. Terms such as allocation (219 occurrences),
investment (272), resource (579), and resources (713) warrant
examination of their operational definitions. Statistical or ecological
compounds matching a seed word can be classification errors rather than
evidence of economic framing.

\subsection{Historical Interpretation}\label{historical-interpretation}

Historical readings of caste and superorganism terminology provide
context
\citep{wheeler1911, gordon1992wittgenstein, boomsma2018superorganismality}.
The BHL era analysis in Section \ref{sec:bhl_grounding} reports literal
OCR frequencies and explicitly sampled stack outputs. Changes in those
rates do not independently establish changes in conceptual commitments,
the origin of a concept, or progress toward more accurate biological
explanation.

\begin{figure}[h]
\centering
\includegraphics[width=0.95\textwidth,height=0.78\textheight,keepaspectratio]{/Volumes/GitVault/Local/worktrees/ento-v1.3.1/ento_linguistics/output/figures/domain_overview_grid.png}
\caption{Ten highest-frequency extracted terms per domain. Bar length is corpus frequency and color is the attached sentence-context entropy estimate. Zero-valued defaults for insufficient contexts must not be interpreted as evidence of unambiguous meaning. These are descriptive extraction outputs from the source-identified abstract layer.}
\label{fig:domain_overview_grid}
\end{figure}

\begin{figure}[h]
\centering
\includegraphics[width=0.95\textwidth,height=0.78\textheight,keepaspectratio]{/Volumes/GitVault/Local/worktrees/ento-v1.3.1/ento_linguistics/output/figures/domain_patterns_grid.png}
\caption{Surface word-formation composition of assigned vocabularies. Each term is assigned to exactly one surface category (single word, multi-word phrase, hyphenated or underscore compound, or containing digits), using the same classifier as Figure \ref{fig:unit_individuality_patterns}. The headline extractor does not independently add n-grams, so availability of a category in the plotting utility does not establish its extraction.}
\label{fig:domain_patterns_grid}
\end{figure}

\begin{figure}[h]
\centering
\includegraphics[width=0.9\textwidth,height=0.78\textheight,keepaspectratio]{/Volumes/GitVault/Local/worktrees/ento-v1.3.1/ento_linguistics/output/figures/unit_of_individuality_patterns.png}
\caption{Unit of Individuality term-formation counts (left) and counts of term names matching scale keyword groups (right). Scale groups can overlap; true zeros are retained. Neither panel estimates actual biological scale boundaries.}
\label{fig:unit_individuality_patterns}
\end{figure}

\begin{figure}[h]
\centering
\includegraphics[width=0.9\textwidth,height=0.78\textheight,keepaspectratio]{/Volumes/GitVault/Local/worktrees/ento-v1.3.1/ento_linguistics/output/figures/concept_hierarchy.png}
\caption{Overlap strength of the six predefined concept categories: each category's weighted degree, the sum of its overlap coefficients with the other categories (Eq.~\ref{eq:overlap_coefficient}), ranked (left) and plotted against associated-term counts (right). Every category links to every other, so unweighted degree would not distinguish them. The display does not represent a biological command hierarchy.}
\label{fig:concept_hierarchy}
\end{figure}

\begin{figure}[h]
\centering
\includegraphics[width=0.9\textwidth,height=0.78\textheight,keepaspectratio]{/Volumes/GitVault/Local/worktrees/ento-v1.3.1/ento_linguistics/output/figures/power_and_labor_term_frequencies.png}
\caption{The fifteen highest-frequency assigned Power and Labor terms in the source-identified abstract layer. Frequency is annotated; color tracks rank. These observations do not quantify hierarchical control or bias.}
\label{fig:power_labor_frequencies}
\end{figure}

\begin{figure}[h]
\centering
\includegraphics[width=0.9\textwidth,height=0.78\textheight,keepaspectratio]{/Volumes/GitVault/Local/worktrees/ento-v1.3.1/ento_linguistics/output/figures/power_and_labor_ambiguities.png}
\caption{Power and Labor terms ranked by attached context-cluster entropy, with corpus frequency and short extraction-context counts. Sentence-context entropy and the plotted extraction-window counts use different context definitions. Cluster entropy is not an independently validated ambiguity measure.}
\label{fig:power_labor_ambiguities}
\end{figure}

\clearpage

\section{Discussion}\label{sec:discussion}

\subsection{Language and Scientific
Practice}\label{language-and-scientific-practice}

The descriptive outputs organize terminology into a framework for
examining the relationship between language and biological explanation.
They are compatible with asking whether scientific metaphors influence
inquiry \citep{latour1987, longino1990, lakoff1980metaphors}, but do not
test that causal hypothesis. Keyword frequency, sentence-cluster
entropy, and co-occurrence identify lexical patterns rather than
researchers' beliefs, decisions, or errors.

Terms such as \emph{queen}, \emph{worker}, and \emph{caste} deserve
contextual scrutiny because biological roles and ordinary-language
connotations can differ. Existing discussions of terminology reform
provide a substantive motivation \citep{herbers2006, herbers2007}. The
present corpus analysis neither establishes that terminology delayed
particular discoveries nor measures the adoption of alternatives across
the field. Such claims would require dated source analysis and evidence
about research decisions.

The observed graph has clustering coefficient 0.8955, while 11.9\% of
assigned terms receive multiple labels. Neither quantity demonstrates
self-reinforcing conceptual bias. Co-occurrence can arise because papers
discuss several biological processes together; label overlap also
follows from the predefined lexicons. A visualization's arrangement must
not be read as a human-style command hierarchy.

\subsection{Network Structure and Its Conditional
Reference}\label{network-structure-and-its-conditional-reference}

A dense co-occurrence graph needs a reference that accounts for
vocabulary selection and document-level opportunity. The companion
fixed-margin extension preserves selected-term counts within documents
and the document frequency of each term. In the executed finite-chain
comparisons, the observed graph has fewer edges and lower clustering
than the randomized reference. The same direction appears under the
longer-burn, wider-spacing sensitivity protocol. The substantive
interpretation is concentration of co-occurrence relative to those
margins, rather than exceptional density attributable to terminology
alone. Topic, genre and other sources of document organization remain
possible explanations. A graph's visual density is therefore an
observation to explain, not evidence of conceptual bias by itself.

The extension also makes vocabulary sensitivity explicit. Changing the
frequency-ranked subset changes the projection and the amount of
connectivity it can display. Interpretation should state the selected
vocabulary, the incidence margins and the comparison model. The
comparison is separately receipted, so its inputs and outputs can be
inspected independently of the corpus analysis.

\subsection{Biological Precision and Terminological
Continuity}\label{biological-precision-and-terminological-continuity}

A terminology decision has several possible costs: obscuring a
biological distinction, inviting an unsupported analogy, or making
relevant literature harder to retrieve. These costs need not move
together. For example, replacing a developmental category with a task
description may lose information even when the replacement sounds less
anthropomorphic. The Introduction's distinction between a definitional
question (what counts as a caste) and a mechanistic one (how experience
shapes division of labor) is the distinction a useful glossary should
preserve.

The practical unit of revision is therefore a defined use in a passage.
Authors can specify the organism or collective being measured, the
criteria for category membership, and the mechanism supported by their
observations. Linking an alternative expression to established indexing
vocabulary can retain discoverability. CACE can structure that review,
but its aggregate should accompany the definition and evidential
rationale rather than decide which expression is biologically correct.

\subsection{Active Inference as a Theoretical
Perspective}\label{active-inference-as-a-theoretical-perspective}

Active Inference offers a vocabulary for discussing generative models,
inference, and action \citep{friston2010free, clark2013whatever}. In the
Active Inferants study, a simulated ant-foraging model reproduces
selected colony phenomena in a laboratory-inspired setting
\citep{friedman2021active}. This provides an example of mechanistic
modeling without a centralized controller; it does not experimentally
compare terminology choices or establish the empirical adequacy of every
biological assumption.

Interpreting terminology as a prior is an analogy here; as the
Introduction states, no Active Inference model is fitted to the corpus.
Whether an ant or colony admits a useful blanket description depends on
specified state variables, dynamics, and conditional-independence
assumptions \citep{friston2013life, kirchhoff2018markov}; it cannot be
settled by changing a noun.

The Environment-Centric Active Inference and related multiscale
proposals in Section \ref{sec:supplemental_analysis} should be read as
theoretical extensions. They require explicit models and empirical tests
before supporting biological or linguistic conclusions.

\subsection{A Complex Systems Reading of
Terminology}\label{a-complex-systems-reading-of-terminology}

The framework can organize questions about scale, connectivity and
feedback without treating a lexical graph as a biological system. For
each passage, specify the units (individuals, tasks, interactions or
colonies), the relevant state variables, the observation timescale and
the coupling being measured. The paper's three networks must remain
distinct: observed term co-occurrence, configured category overlap and
the biological interactions discussed by the underlying literature. A
dense term graph does not supply evidence that the ants have an equally
dense communication graph.

This distinction supports productive comparison rather than automatic
analogy. A mechanism-based account of task allocation can connect the
Power and Labor vocabulary to questions about local information and
regulation, while biological individuality asks whether the reported
outcome is measured at individual or colony scale. The next empirical
step is to annotate these mechanistic commitments in source passages and
test agreement on held-out documents. Only then can the framework assess
whether a terminology choice helps readers recover the actual units,
coupling and conditions of a study. Expanding to other taxa, genres or
languages requires a new measurement-validation assessment; reproducible
software alone does not establish transportability.

\subsection{Practical Use of CACE}\label{practical-use-of-cace}

CACE makes evaluation criteria inspectable: clarity of operational
definitions, suitability of metaphors, consistency of use, and
adaptability to new findings. Its numerical implementation is a
heuristic instantiation. Penalizing membership in a predefined
anthropomorphic vocabulary partly encodes the desired ranking; a
favorable score is not independent validation of a replacement term.

The representative-term table includes terms absent from the extraction.
Their scores are derived from text features and fallback conventions, as
recorded in the Extracted column. In particular, a zero default entropy
for an unattested term must not be interpreted as demonstrated
specificity. Comparisons between \emph{slave} and \emph{host worker}
illustrate these rules, rather than measuring a change in readers'
comprehension or research outcomes.

The computed aggregate values are 0.30 and 0.47, respectively, with
Appropriateness 0.50 and 0.40. These values remain useful for inspecting
the scoring implementation. Independent assessment should compare
definitions and context-specific biological accuracy, collect blinded
judgments, measure agreement, and test sensitivity to the chosen weights
before treating the scores as a prescriptive standard.

\subsection{Cross-Domain
Communication}\label{cross-domain-communication}

Domain overlap can help select terms for shared glossaries and explicit
operational definitions. The six-domain framework is one proposed
partition; alternative lexicons and annotation schemes may produce
different assignments. Maintaining links to established terminology can
support discoverability while clarifying which mechanism or
observational category is meant. Benefits to communication remain
hypotheses for reader and author studies.

\subsection{Measurement Validation and Terminology
Intervention}\label{measurement-validation-and-terminology-intervention}

Two study designs separate measurement quality from communication
effects. The first is the annotation design specified in Section
\ref{sec:measurement_validation}: it tests whether occupancy entropy and
lexical indicators track reader-labeled meanings and discourse functions
after accounting for term frequency, document length and topic, with
agreement and uncertainty reported at the level of the annotated
observation.

The second borrows the manipulation logic of metaphor experiments
(Section \ref{sec:related_work}). A randomized reader study can hold
biological evidence constant while varying an established label, a
proposed alternative and an explicit operational definition.
Prespecified outcomes include inference accuracy, confidence
calibration, recall and literature retrieval. This design tests a
terminology effect directly. Comparing its outcomes with CACE rankings
then evaluates the scoring proposal using evidence outside the scheme.
Researchers and students can be included as prespecified groups rather
than assumed to interpret a term identically.

\subsection{Limitations}\label{limitations}

\begin{enumerate}
\def\labelenumi{\arabic{enumi}.}
\tightlist
\item
  \textbf{Source custody and relevance.} Some archived abstract strings
  lack digest-indexed source metadata and are excluded from headline
  analysis; broad queries include adjacent biology and computational
  uses. Historical volumes contain mixed topics; PMC notices are
  excluded by a title rule rather than a complete article-type screen.
  Source reconciliation and relevance annotation remain incomplete.
\item
  \textbf{Sampling and accessibility.} Search and availability filters
  select a convenience corpus. Layers are not matched document samples
  and differ in extraction threshold, length, genre, and publication
  history.
\item
  \textbf{OCR and language.} Literal historical matching is sensitive to
  OCR errors, language, hyphenation, and spelling. Rare or absent
  matches do not date a concept's origin.
\item
  \textbf{Proxy validity.} Computational clusters are not annotated
  senses; marker patterns are not validated discourse labels. The
  pipeline does not measure author intent, cognitive distortion, or
  causal framing.
\item
  \textbf{Statistical dependence.} Domain groups overlap and share
  document-derived contexts. Reported test outputs are exploratory; no
  calibrated confirmatory inference is claimed.
\item
  \textbf{Explicitly bounded components.} PMC discourse uses an
  approximately one-fifth deterministic text sample; domain CACE uses at
  most fifty terms; BHL entropy uses twenty frequent terms per era with
  full-corpus contexts. Default BHL extraction, framing, and
  literal-frequency calculations include all stored documents; an
  optional development character cap produces a separately fingerprinted
  subset.
\item
  \textbf{Unmeasured theory.} No fitted generative model, empirical
  Markov-blanket estimation, term-reform intervention, or independent
  human validation is reported.
\end{enumerate}

These boundaries support a focused next study: a frozen, fully
reconciled corpus with independent annotations and prespecified
document-level analysis, followed by a controlled evaluation of
operational terminology.

\subsection{Complexity as an Explicit Modeling
Question}\label{complexity-as-an-explicit-modeling-question}

Complexity connects the paper's questions about organization,
information and scale while requiring a defined target system and
measurement. Accounts of complex systems distinguish candidate features
and mathematical measures rather than supplying a universal
interpretation of density or entropy \citep{ladyman2013complex}. A
complete graph has maximal density and a short description; dense word
co-occurrence alone does not establish emergence, adaptation or
collective cognition. One-mode projection itself can create local
clustering \citep{newman2003structure}. The new representation
diagnostic consequently reports how graph statistics depend on
vocabulary and edge thresholds, without converting them into biological
properties.

Architecture asks which units and relations a passage specifies. Simon's
near-decomposability account concerns relatively weak coupling between
subsystems and different short- and long-run behavior
\citep{simon1962architecture}; overlapping lexical domains do not
demonstrate those dynamical conditions. Fertilization exchanges
questions and models across collective behavior, network science and
terminology research while retaining their observational units.
Expansion asks whether measurement transfers to new taxa, genres,
languages and readers after held-out validation. These three themes
support a complex-systems research program; the executed corpus analyses
do not by themselves complete it.

\clearpage

\section{Conclusion}\label{sec:conclusion}

This work provides a six-domain framework and a reproducible descriptive
pipeline for examining terminology in scientific text. Across 7599
source-identified abstracts, the pipeline processes 998,645 tokens and
extracts 11,700 candidate terms, with 1330 receiving domain assignments.
Observed document co-occurrence, predefined conceptual-category overlap,
sentence-context entropy, and heuristic framing and CACE scores are
reported as distinct quantities.

The 11.9\% of assigned terms with multiple labels measures
classification overlap, not semantic drift. The 6 concept categories are
predefined rather than discovered. Complementary source layers extend
the descriptive scope without establishing that language causes bias,
that a terminology reform improves scientific modeling, or that
numerical CACE rankings are independently validated.

The fixed-margin comparison sharpens the network contribution: visual
density must be interpreted relative to document and term incidence, and
the executed reference indicates more concentrated co-occurrence than
those margins alone would predict. It supplies a conditional descriptive
comparison without identifying a linguistic cause. The broader
contribution is a method for keeping biological definitions, lexical
observations, and proposed communication effects distinguishable while
investigating their relationships.

\subsection{Future Directions}\label{future-directions}

The immediate research priorities are source reconciliation,
document-level relevance and license review, and independent annotation
of senses and framing. A frozen annotated corpus would support tests of
extraction accuracy, agreement between annotators, context-count
sensitivity, and statistical models accounting for shared documents and
overlapping domain labels.

Longitudinal analysis should distinguish changes in source composition
and OCR quality from changes in terminology. Multilingual comparisons
require language-specific normalization and definitions. Reader or
author experiments could then evaluate whether operational definitions
and alternative terms improve comprehension or mechanistic explanation.

CACE offers explicit evaluation questions for such studies. Its current
scores provide an inspectable starting point, with design choices and
missing-data conventions disclosed. Active Inference and multiscale
interpretations remain theoretical proposals until connected to
specified models and empirical measurements. The repository contributes
tools and auditable descriptive results for that work.

The complexity theme connects precise units, cross-disciplinary exchange
and validated transfer. Architecture specifies what is represented;
Fertilization exchanges questions while preserving evidence boundaries;
Expansion tests measurement in new settings. The vocabulary--threshold
diagnostic and graphical abstract make that program explicit.
Reproducible representation is an executed contribution, whereas
annotated meaning, reader effects and biological dynamics remain
separate validation targets.

\clearpage

\section{Related Work}\label{sec:related_work}

\subsection{Scientific Language and
Categories}\label{scientific-language-and-categories}

Philosophy of science and discourse research provide context for
studying terminology as part of research practice
\citep{kuhn1996, latour1987, longino1990, fairclough1992, wodak2009methods}.
Conceptual and deliberate-metaphor frameworks supply complementary
questions about ordinary connotations and communicative intention
\citep{lakoff1980metaphors, steen2017deliberate}. This study draws
motivation from those traditions, while keeping computational lexical
proxies separate from claims about thought, intention, or causal
influence.

Work on scientific classification and cultural histories offers
additional context
\citep{berlin1992, hacking1999social, sleigh2007ants}. The present
six-domain taxonomy is a proposed analytic organization; it is not
independently established as exhaustive or uniquely appropriate.

\subsection{Terminology in Social-Insect
Research}\label{terminology-in-social-insect-research}

Debates over caste categories and descriptions of ant behavior motivate
explicit definitions
\citep{gordon1992wittgenstein, boomsma2018superorganismality}. Research
on collective behavior provides biological context for distinguishing
task allocation from an assumed centralized controller
\citep{gordon2010, gordon2019ecology, gordon2023ecology}. The
descriptive corpus pipeline does not itself test those mechanisms.

Discussions of racially loaded language in social-insect research
provide a substantive case for examining terminology choices
\citep{herbers2006, herbers2007}. The ESA Better Common Names Project is
a related institutional initiative \citep{betternamesproject2024}, and
the term \emph{colony} itself has been examined for its settler-colonial
connotations \citep{vis2026colony}. Neither the current frequency
analysis nor CACE scoring measures the adoption or effects of these
reforms.

Molecular accounts of caste and epigenetic mechanisms
\citep{sumner2018molecular, oldroyd2021epigenetics} and developmental
canalization research \citep{qiu2022canalized} reinforce the need to
distinguish developmental processes from temporary behavioral labels.
Canalization must not be interpreted as proof that every caste category
is labile. The historical superorganism literature \citep{wheeler1911}
and later conceptual analysis \citep{boomsma2018superorganismality} also
caution against dating a concept's origin from an absent OCR match.

\subsection{Computational Mapping and Source-Layer
Analysis}\label{computational-mapping-and-source-layer-analysis}

Computational literature mapping supplies methods for examining
connections among terms and publications \citep{chen2006citespace}.
Here, document co-occurrence is kept separate from vocabulary-overlap
relationships between predefined concepts. TF-IDF/KMeans occupancy
entropy is used as a descriptive measure of context distributions,
without claiming independent sense annotation or validated linguistic
ambiguity.

Source-layer comparisons are also descriptive. Abstracts, full texts,
historical volumes and preprints differ in access, genre, length, topic
and extraction threshold. An observed difference between layers is not
automatically a robustness result or a historical change in meaning.

\subsection{Validation and Experimental
Evidence}\label{validation-and-experimental-evidence}

\citet{grimmer2013text} emphasize, for automated content analysis,
validation tailored to the substantive question. Its relevance here is
methodological: word counts and computational categories require an
explicit connection to the construct they are intended to measure. This
study's receipts establish computational custody, not construct
validity; Section \ref{sec:measurement_validation} specifies the
annotation design that would address the latter.

Experimental metaphor research supplies a complementary way to
investigate reader responses. \citet{thibodeau2011metaphors} studied how
contrasting crime metaphors affected reasoning. That evidence concerns a
different topic and participant setting. It supports using controlled
language manipulations to formulate testable questions, while leaving
the effects of technical ant terminology on specialists and students
unresolved. Connecting these traditions requires both measurement
validation and a direct communication experiment.

\subsection{Active Inference and Colony
Modeling}\label{active-inference-and-colony-modeling}

The Free Energy Principle and Active Inference supply theoretical
vocabulary for generative modeling
\citep{friston2010free, friston2013life, clark2013whatever, kirchhoff2018markov}.
The Active Inferants framework supplies a simulated ant-foraging example
\citep{friedman2021active}; model behavior alone cannot establish how
terminology choices affect scientific reasoning.

Theoretical perspectives on eusociality and biological individuality
\citep{nowak2010evolution, boomsma2018superorganismality} motivate
careful definition of units and mechanisms. The Environment-Centric
Active Inference extensions in the supplement remain proposed
constructions requiring explicit models and empirical testing.

\subsection{Collective Behavior and Complex
Systems}\label{collective-behavior-and-complex-systems}

Collective-behavior research supplies a mechanistic counterpart to the
terminology framework. \citet{couzin2009collective} relates group-level
response to social interaction, individual state and environmental
modification. \citet{feinerman2017cognition} distinguishes collective
responses supported by individual information from outcomes that require
interactions across a scale gap. These accounts motivate asking which
entity holds information and which process combines it; a colony-level
description alone does not identify the computation.

Ecological conditions also matter to how interaction processes regulate
activity \citep{gordon2014ecology}. In an experimental harvester-ant
example, combined food and forager chemical cues on mimics increased
outgoing foraging activity \citep{greene2013chemical}. Such experiments
link a defined intervention to an observed response. They illustrate the
level of biological evidence needed for a mechanistic claim, rather than
validating any particular lexical label. The corpus network instead
records pairs of expressions in documents: its vertices, edges and
timescale differ from those of a measured ant-interaction network.
Similar graph vocabulary should prompt an explicit comparison of
observational units before a mechanism is transferred between fields.

\subsection{Positioning This Work}\label{positioning-this-work}

This repository contributes a six-domain descriptive workflow with
auditable source layers, inspectable computational definitions,
regenerable figures and receipt-bound values. CACE is a proposed
evaluation framework, not a validated measure or intervention. Its
independent validation would require annotated meanings, blinded
judgments, measured agreement, and tests of sensitivity and
communication outcomes. The distinction between implemented measurement,
theoretical motivation, and unperformed validation is part of the
contribution.

\subsection{Complexity, Representations, and
Dynamics}\label{complexity-representations-and-dynamics}

Simon \citep{simon1962architecture} motivates explicit units, coupling
and timescales in hierarchical systems. Newman
\citep{newman2003structure} distinguishes network structure, models and
dynamical processes, including clustering induced by affiliation
projection. Ladyman and colleagues \citep{ladyman2013complex} compare
candidate features and different measures of complexity. Together these
sources support precise modeling questions rather than equating a
lexical graph with the biological system it describes. The present
contribution is an auditable discourse representation and its
sensitivity analysis; testing biological dynamics or readers' inferences
requires separate observations and designs.

\clearpage

\section{Acknowledgments}\label{sec:acknowledgments}

We gratefully acknowledge the contributions of individuals and
institutions that made this research possible.

\subsection{Institutional Support}\label{institutional-support}

This work was conducted at the Active Inference Institute and APOIDEAS.
We thank the Institute for providing the research environment and
collaborative infrastructure that supported the development of the
Ento-Linguistic framework.

\subsection{Collaborations}\label{collaborations}

We thank colleagues and collaborators for valuable discussions and
feedback throughout the development of this work, particularly regarding
the six-domain framework and the design of the descriptive pipeline.

\subsection{Data and Software}\label{data-and-software}

This research builds upon open-source software tools and publicly
available datasets. We acknowledge:

\begin{itemize}
\tightlist
\item
  Python scientific computing stack (NumPy, SciPy, Matplotlib, NetworkX)
\item
  Natural Language Toolkit (NLTK) for text processing and scikit-learn
  for TF-IDF vectorization and clustering
\item
  LaTeX and Pandoc for document preparation
\item
  Published entomological literature informing the domain terminology
  seeds
\end{itemize}

\begin{center}\rule{0.5\linewidth}{0.5pt}\end{center}

\emph{All errors and omissions remain the sole responsibility of the
authors.}

\clearpage

\section{Supplemental Methods: Text Processing and Term
Extraction}\label{sec:supplemental_methods}

This section specifies the input and extraction stages used by the
study. The numerical definitions and bounded components are detailed in
Section \ref{sec:supplemental_infrastructure}. Developer APIs remain
documented in source; support utilities are not additional research
measurements.

\subsection{Source Layers and Custody}\label{source-layers-and-custody}

The archived abstract input contains 7609 strings. The headline analysis
includes 7599 strings whose SHA-256 digest maps to an identified PubMed
record and excludes 10 unreconciled strings. Exact digest matches to
retrieved PubMed abstracts can recover metadata without changing
archived text. Identification establishes a source association, not
relevance to ant biology, a complete retrieval history, or independent
text annotation.

PMC records retain title, abstract, and body text in their own shard
directory. BHL records retain historical OCR and era assignments. arXiv
title/abstract records form a separate preprint layer. Those documents
are not concatenated into the headline abstract input. The custody audit
inspects every stored record and reports missing metadata, duplicate
text, identifier disagreement, unused sidecar entries, and invalid text.

Digest-indexed sidecars can retain only one source identity when several
records share identical text. These collisions are reported rather than
resolved by inventing metadata. Broad searches include adjacent biology,
computational terminology and mixed-topic historical volumes. The stored
PMC layer includes correction and retraction notices whose repeated
boilerplate produces shared body text. Records whose titles begin with a
correction, erratum or retraction prefix, including retracted full
articles, are retained in the shards but excluded from full-text
analysis. Individual relevance, article-type screening and license
review remain incomplete.

\subsection{Normalization and Token
Streams}\label{normalization-and-token-streams}

TextProcessor applies Unicode normalization, lowercase conversion, word
tokenization, punctuation filtering, and stop-word removal. Hyphenated
and underscored tokens can remain intact. The stop vocabulary combines
NLTK English stop words with configured scientific meta-language words.
WordNet lemmatization is available and is used for corpus vocabulary
summaries; terminology extraction disables lemmatization for its
occurrence counts. Consequently a summary's most-common lemma and an
extracted surface-form frequency need not have the same spelling.

NLTK sentence tokenization supplies the contexts used for semantic
entropy. These sentence contexts differ from the short, processed-token
windows retained on extracted Term objects. English-language resources
are applied to a heterogeneous stored corpus; this is not a validated
multilingual processing pipeline.

NLTK resources are separate downloaded prerequisites. The analysis
receipt binds the selected English stopword file, English punkt\_tab
files and WordNet dictionary bytes (the complete archive when
archive-backed). The same hashes enter cache signatures; changing a
selected resource requires reanalysis. These hashes identify the
installed inputs without vendoring them or guaranteeing automatic
restoration. Open Multilingual WordNet is outside this receipt because
English lemmatization does not consume it. The source implementation and
uv.lock specify Python package dependencies.

\subsection{Candidate Extraction and Domain
Assignment}\label{candidate-extraction-and-domain-assignment}

TerminologyExtractor counts tokens, applies the layer-specific minimum
frequency listed in Methods, and evaluates candidate filters.

Candidate filters retain configured seed words of three to fifty
characters before evaluating generic filters. Other candidates enter
through scientific patterns, compound separators, or configured
scientific substrings. They reject pure numbers but can admit unrelated
substring matches and OCR artifacts. Candidate inclusion is therefore
broader than biological terminology. For example, skin, making and
queensland can remain candidates without receiving a domain label.

The six DOMAIN\_SEEDS vocabularies define direct assignments. Compound
tokens can inherit labels from seed-word overlap; fallback lexical
patterns use word boundaries. An extracted term can receive several
labels. Label overlap follows from these rules and does not demonstrate
temporal semantic drift or independently annotated meanings.

Candidate iteration is sorted for deterministic output. The extractor
stores surface text, lemma, frequency, domain labels, confidence, and
deduplicated short contexts.

Extraction confidence combines configured frequency, context and
classification features. It is not calibrated against human correctness
labels. The pipeline does not report precision, recall, multilingual
accuracy, or inter-rater reliability.

\subsection{Document Co-occurrence and Concept
Categories}\label{document-co-occurrence-and-concept-categories}

The observed terminology graph uses the hundred highest-frequency
domain-assigned terms, with lexical tie-breaking. Case-insensitive
whole-word searches identify presence in each document. A pair's edge
weight counts documents containing both terms, regardless of repeated
mentions. Shared labels do not create observed edges. The saved graph
statistics describe this selected vocabulary rather than the entire
field.

The separate concept map uses six predefined categories and
vocabulary-association rules. Terms can belong to several categories, so
a total of term associations is not a count of unique terms. Category
links summarize vocabulary overlap. Node positions, direct-link
centrality, colors, and label selection are display conventions and do
not establish biological hierarchy or causal dependence.

The overlap heatmap uses the Szymkiewicz--Simpson coefficient in
Equation \ref{eq:overlap_coefficient}. Its diagonal is one by
construction. A zero observed overlap does not establish conceptual
separation.

\subsection{Reproducible Inputs and Failure
Behavior}\label{reproducible-inputs-and-failure-behavior}

DataLoader rejects malformed or empty corpus input, including non-string
and blank records. Required pipeline stages propagate failures. Process
workers merge results in input order; ENTO\_ANALYSIS\_WORKERS=1 selects
serial execution. Resource-related pool unavailability is logged before
a serial retry, while other worker failures propagate.

Caches bind ordered record contents, project Python source, locked
dependencies, and explicit development limits. Changed body text
invalidates an artifact even if its metadata, record count and file
modification time remain unchanged. The completed receipt binds the
input/output inventory and image hashes. A receipt establishes which
bytes were used, not the scientific validity of the taxonomy or proxies.

\clearpage

\section{Supplemental Methods: Statistical and Scoring
Infrastructure}\label{sec:supplemental_infrastructure}

The canonical implementations are semantic\_entropy, DomainAnalyzer,
statistics\_pipeline, fulltext\_pipeline, arxiv\_analysis, and
bhl\_analysis. Their generated artifacts preserve exclusions, coverage,
and exploratory numerical outputs.

\subsection{Sentence-context Entropy}\label{sentence-context-entropy}

Each source text is sentence-tokenized once for a term set. An inverted
index over alphanumeric parts prunes impossible matches; a
case-insensitive whole-word regular expression makes the final decision.
This preserves matching sentence order. Contexts must contain more than
three words, and at least five usable contexts are required for
computation.

TF-IDF uses English stop words, minimum document frequency one, and at
most one thousand features. KMeans uses random state 42 and ten
initializations. The requested cluster count follows the bounded
square-root rule in Methods, remaining below the number of contexts.
Occupied clusters may be fewer than requested.

Shannon entropy is computed from occupied-cluster probabilities in bits
(Equation \ref{eq:semantic_entropy}). The normalized value divides by
the occupied-cluster ceiling; a single occupied cluster has zero
normalized entropy. The exploratory high-entropy flag uses a threshold
above two bits. Computational partitions are not independently annotated
word senses.

Results record ok, insufficient\_contexts, or error status. Only ok
values enter domain means and statistical groups. Excluded estimates are
counted separately; their placeholder zero is not included as a measured
value. A term assigned to several domains contributes to each domain
group. Entropy-table sample sizes and weighted totals count valid domain
memberships, not unique terms or documents.

\subsection{Exploratory Comparisons and
Intervals}\label{exploratory-comparisons-and-intervals}

The statistical artifact reports per-domain valid-entropy means,
standard deviations where at least two values exist, and exclusions.
Welch tests compare groups with at least two values. Benjamini--Hochberg
adjustments apply to the emitted comparison family at a nominal
threshold of 0.05. Standardized differences use a small-sample bias
correction; they should not be described as uncorrected Cohen's d.

One-way ANOVA and its effect-size summary use the valid groups. Figure
intervals use nominal Student-t quantiles and are omitted for groups
with only one valid estimate. These numerical calculations have software
tests and independent numerical comparators, but their simple population
interpretation is not justified by this corpus design.

Terms can share documents, contexts and domain labels. Multiplicity
adjustment does not repair that dependence, convenience retrieval,
unequal context counts, or differing extraction thresholds. Threshold
flags are exploratory. Confirmatory inference requires a prespecified
document-level model and independently annotated evaluation data.

\subsection{Shared CACE Sample and
Scores}\label{shared-cace-sample-and-scores}

Domain figures and statistical exports use the same scoring function and
sample: up to fifty terms per domain, ordered by decreasing frequency
with lexical tie-breaking. Each term supplies at most ten of its stored
short extraction contexts. The representative-term table is a separate
named-term evaluation whose artifact records whether a term was
extracted.

Clarity uses attached sentence-context entropy when available.
Appropriateness uses configured vocabulary and domain-overlap penalties.
Consistency uses TF-IDF context-vector cosine similarities. Evolvability
uses assigned domains and scale-marker categories. Their equally
weighted mean is the aggregate; Equations
\ref{eq:cace_clarity}--\ref{eq:cace_evolvability} specify the proposed
rules.

Unavailable entropy still defaults to zero inside the heuristic, and
fewer than two contexts yield a Consistency convention of 0.5. These
conventions do not establish clarity or consistency. Missing figure
input is labeled unavailable instead of substituting confidence as a
CACE score. No completed coefficient-sensitivity study, human-rating
validation, intervention benefit, or calibrated scientific-accuracy
prediction is claimed.

\subsection{Framing and Discourse
Proxies}\label{framing-and-discourse-proxies}

Occurrence-context framing uses three-token windows around
domain-assigned terms. The feature extractor has four anthropomorphic,
four hierarchical, and four economic regex patterns. Marker matches are
lexical events, not validated author intent or judgments that a
biological description is misleading.

Domain framing proportions count occurrence contexts matching an
anthropomorphic pattern. Multi-domain occurrences contribute to each
domain but once to the overall tally. Distinct framing-marked term
counts in the figure are a different summary from those occurrence
proportions.

Discourse, rhetorical, and persuasive analyzers use configured patterns
and keywords. Their artifacts record eligible/analyzed texts and
exclusions. PMC discourse uses a deterministic stride sample when corpus
size exceeds the configured text/character targets, with an
approximately one-fifth sample in this snapshot. Raw counts depend on
text length and sample size; the cross-layer figure is not a normalized
prevalence comparison or an annotated discourse study.

\subsection{Historical OCR and Layer
Coverage}\label{historical-ocr-and-layer-coverage}

BHL literal frequencies normalize NFC text, lowercase it, join
hyphenated line breaks, and tokenize Latin-letter words while retaining
internal hyphens. Seed phrases match exact consecutive tokens; counts
are non-overlapping per phrase. Rates divide by era token counts and
multiply by ten thousand. Alternative spellings, OCR errors and language
affect those rates. An observed zero does not date a concept's origin.

The BHL computational stack uses cleaned OCR with the canonical
extraction and classification rules. Default extraction and framing
process all stored era documents using per-document counters, avoiding a
corpus-wide token-position map. Stored and cleaned text are still held
in memory; memory use is not constant in corpus size. Entropy evaluates
twenty frequent candidates per era with sentence contexts drawn from all
era documents. Unassigned candidates can enter the overall sampled
entropy mean, while domain means only use assigned valid candidates.

Coverage records available and analyzed documents and characters. An
explicit positive \texttt{BHL\_STACK\_CHARACTER\_BUDGET} selects a
deterministic whole-document development subset.
\texttt{FULLTEXT\_ANALYSIS\_LIMIT} similarly bounds a development PMC
run. Both limits enter fingerprints; bounded caches cannot satisfy an
unbounded default run.

\subsection{Execution and Artifact
Verification}\label{execution-and-artifact-verification}

The standalone sequence installs locked dependencies and NLTK resources,
runs tests, generates figures, checks custody, and renders the
manuscript. Testing and generation run sequentially because some
integration tests consume generated exports.

The analysis signature also hashes the selected English NLTK tokenizer,
stopwords and WordNet contents. The receipt records these resource
hashes without machine-specific installation paths. Missing resources
fail, and changed resources invalidate caches. These content hashes do
not vendor or automatically restore external resource downloads.

Reported values are injected from receipted artifacts at render time.
The generator checks finite/nonempty JSON, decodable PNGs, registry
hashes, figure inventory and manuscript figure availability before
writing its receipt. Rendering validates the receipt, rejects unresolved
placeholders, and requires successful Pandoc, XeLaTeX and BibTeX
execution. Undefined citations/references and missing glyphs fail the
final build.

Software tests use actual corpus slices, numerical examples, real files,
local HTTP and real subprocesses. Negative controls establish that
malformed input, stale artifacts and failing commands are rejected.
Those checks establish software behavior and provenance, not source
relevance, individual reuse rights, human-validation outcomes or causal
scientific conclusions.

Completed BHL eras are saved atomically in local recovery checkpoints.
Each checkpoint binds the ordered era records,
implementation/dependency/resource signature and any development bound,
plus a digest of the completed result. A matching checkpoint can resume
a disrupted run; changed inputs or bounds require recomputation, and
corrupted matching results fail. The final four-layer receipt is still
written only after all required stages complete.

\subsection{Fixed-Margin Network Robustness
Extension}\label{sec:fixed_margin}

A separately receipted companion workflow reconstructs the abstract
terminology network from the frozen identified records and the same
frequency-ranked vocabulary. It randomizes the binary document--term
incidence matrix with Curveball trades
\citep{strona2014curveball, carstens2018curveball}. A trade retains
common terms in both selected documents and uniformly repartitions their
exclusive terms while preserving both row sizes. Consequently, each term
retains its document frequency and each document retains its number of
selected terms. Self-loop and other no-op transitions remain part of the
chain. The total weight of the projected graph equals the sum of the
number of term pairs in each document; this quantity is an invariant of
the conditioning margins.

Every retained draw checks the two margins and the weight invariant. The
workflow reports projected edge count, mean local unweighted clustering,
retained-chain traces, between-chain diagnostics and empirical
conditional distributions. It also evaluates observed projections under
smaller vocabulary subsets of the same ranking. A second protocol uses
longer burn-in and spacing with different seeds. These checks assess
structural interpretation and sampling sensitivity. Finite correlated
draws and between-chain agreement do not prove convergence to the
stationary distribution. Conditional null envelopes describe the sampled
reference and are distinct from population confidence intervals.

The companion implementation is in research/network\_robustness/. Its
receipt binds the published-core input receipt, exact source inputs,
extension implementation, dependency lock and generated outputs. It does
not relabel cached core results as newly recomputed analyses. The
executed report and trace array are provided as companion artifacts.

A validation guard additionally decodes retained numerical traces and
the companion PNG, verifies full protocol dimensions, finite values and
graph bounds, reconstructs the observed graph from source-identified
abstracts, and recomputes the JSON summary statistics from all retained
draws. The companion report's result table must agree with those
summaries, and the values reported in the Results section are injected
from the same summaries. Negative controls include self-hashed corrupt
files and inconsistent summaries. These internal-consistency checks do
not establish chain mixing, validate lexical proxies, or certify the
visual interpretation of a decoded figure.

\clearpage

\section{Supplemental Results}\label{sec:supplemental_results}

\subsection{Exploratory Pairwise Domain
Comparisons}\label{exploratory-pairwise-domain-comparisons}

Table \ref{tab:pairwise_domain} presents pairwise comparisons of
per-term semantic entropy between all Ento-Linguistic domains using
Welch's two-sample \(t\)-tests. Raw \(p\)-values are computed from the
\(t\)-distribution with Satterthwaite-approximated degrees of freedom;
adjusted \(p\)-values correct for 15 simultaneous comparisons using the
Benjamini-Hochberg (BH) procedure at \(q = 0.05\). The effect size \(d\)
is a small-sample bias-corrected (Hedges-type) standardized difference.
Conventional benchmarks of about 0.2, 0.5 and 0.8
\citep{cohen1988statistical} are rules of thumb that are not calibrated
for these overlapping, dependent groups. Domain descriptives entering
these tests are per-term valid-entropy means with exclusions counted.

\begin{table}[h]
\centering
\small
\begin{tabular}{|l|l|c|c|c|c|c|}
\hline
\textbf{Domain A} & \textbf{Domain B} & \textbf{$t$} & \textbf{$p$ (raw)} & \textbf{$p$ (BH)} & \textbf{Std.\ diff.\ $d$} & \textbf{Sig.\ (BH)} \\
\hline
Behavior \& Identity & Economics & -0.9801 & 0.3312 & 0.4516 & -0.2166 & no \\
Behavior \& Identity & Kin \& Relatedness & -0.1601 & 0.8736 & 0.8736 & -0.0402 & no \\
Behavior \& Identity & Power \& Labor & 1.2093 & 0.2289 & 0.4516 & 0.2101 & no \\
Behavior \& Identity & Sex \& Reproduction & 0.1864 & 0.8525 & 0.8736 & 0.0337 & no \\
Behavior \& Identity & Unit of Individuality & 2.3743 & 0.0192 & 0.1442 & 0.3793 & no \\
Economics & Kin \& Relatedness & 0.6153 & 0.5417 & 0.6772 & 0.1781 & no \\
Economics & Power \& Labor & 2.0805 & 0.0423 & 0.1586 & 0.4216 & no \\
Economics & Sex \& Reproduction & 1.1366 & 0.2603 & 0.4516 & 0.2350 & no \\
Economics & Unit of Individuality & 3.1260 & 0.0031 & 0.0464 & 0.5875 & yes \\
Kin \& Relatedness & Power \& Labor & 1.0421 & 0.3045 & 0.4516 & 0.2457 & no \\
Kin \& Relatedness & Sex \& Reproduction & 0.2982 & 0.7672 & 0.8736 & 0.0704 & no \\
Kin \& Relatedness & Unit of Individuality & 1.8320 & 0.0767 & 0.2300 & 0.4127 & no \\
Power \& Labor & Sex \& Reproduction & -0.9976 & 0.3203 & 0.4516 & -0.1689 & no \\
Power \& Labor & Unit of Individuality & 1.1416 & 0.2553 & 0.4516 & 0.1687 & no \\
Sex \& Reproduction & Unit of Individuality & 2.1273 & 0.0353 & 0.1586 & 0.3348 & no \\
\hline
\end{tabular}
\caption{Exploratory pairwise Welch tests on valid per-term sentence-context entropy, with Benjamini--Hochberg adjusted $p$-values over 15 comparisons and standardized effect sizes. Threshold flags use $q=0.05$; domain groups share terms and contexts, so the flags are exploratory (Supplemental Section \ref{sec:supplemental_infrastructure}). The omnibus ANOVA yields $F(5,350)=2.4761$, $p$-value 0.0319, and $\eta^2=0.0342$.}
\label{tab:pairwise_domain}
\end{table}

\subsection{CACE Scoring for Key
Terms}\label{cace-scoring-for-key-terms}

Table \ref{tab:cace_full} presents full CACE evaluations for a
representative set of entomological terms, comparing anthropomorphic
labels with proposed functional alternatives.

\begin{table}[h]
\centering
\small
\begin{tabular}{|l|c|c|c|c|c|c|}
\hline
\textbf{Term} & \textbf{C} & \textbf{A} & \textbf{Cs} & \textbf{E} & \textbf{Mean} & \textbf{Extracted} \\
\hline
queen & 0.48 & 0.45 & 0.04 & 0.33 & 0.33 & yes \\
\textit{primary reproductive} & 1.00 & 1.00 & 0.50 & 0.00 & 0.62 & no \\
\hline
worker & 0.42 & 0.45 & 0.05 & 0.33 & 0.31 & yes \\
\textit{non-reproductive helper} & 1.00 & 1.00 & 0.50 & 0.00 & 0.62 & no \\
\hline
slave & 0.31 & 0.50 & 0.05 & 0.33 & 0.30 & yes \\
\textit{host worker} & 1.00 & 0.40 & 0.50 & 0.00 & 0.47 & no \\
\hline
caste & 0.36 & 1.00 & 0.06 & 0.33 & 0.44 & yes \\
\textit{task group} & 1.00 & 1.00 & 0.50 & 0.00 & 0.62 & no \\
\hline
soldier & 0.35 & 0.50 & 0.10 & 0.17 & 0.28 & yes \\
\textit{major worker} & 1.00 & 0.50 & 0.50 & 0.00 & 0.50 & no \\
\hline
colony & 0.47 & 1.00 & 0.03 & 0.33 & 0.46 & yes \\
haplodiploidy & 0.31 & 1.00 & 0.05 & 0.17 & 0.38 & yes \\
trophallaxis & 0.33 & 1.00 & 0.06 & 0.17 & 0.39 & yes \\
\hline
\end{tabular}
\caption{Heuristic representative-term CACE scores: C is Clarity, A Appropriateness, Cs Consistency, and E Evolvability; Mean weights them equally. Extracted indicates membership in the extracted vocabulary, not phrase presence anywhere in the source text. Unextracted alternatives use fallback conventions. Defaults and vocabulary penalties are not independently measured clarity, biological accuracy, or evidence of terminology-reform benefit.}
\label{tab:cace_full}
\end{table}

\subsection{Semantic Entropy
Distribution}\label{semantic-entropy-distribution}

Table \ref{tab:entropy_distribution} summarizes the distribution of
semantic entropy across domains.

\begin{table}[h]
\centering
\begin{tabular}{|l|c|c|c|}
\hline
\textbf{Domain} & \textbf{Mean $H$ (bits)} & \textbf{High-entropy terms (\%)} & \textbf{$N$} \\
\hline
Economics & 1.85 & 46.2 & 26 \\
Power \& Labor & 1.64 & 31.6 & 76 \\
Behavior \& Identity & 1.74 & 44.6 & 56 \\
Sex \& Reproduction & 1.73 & 43.8 & 64 \\
Unit of Individuality & 1.55 & 26.8 & 112 \\
Kin \& Relatedness & 1.76 & 45.5 & 22 \\
\hline
\textbf{Overall} & 1.67 & 36.2 & \textbf{ 356 } \\
\hline
\end{tabular}
\caption{Distribution of sentence-context entropy $H(t)$ across Ento-Linguistic domains (Eq.~\ref{eq:semantic_entropy}; cluster-count rule in Methods). High-entropy terms exceed $H > 2.0$ bits; the statistic summarizes occupancy of computational clusters, not annotated senses. Normalized entropy divides by $\log_2 k_{\mathrm{occupied}}$ and is zero when one cluster is occupied. $N$ counts terms with usable estimates; Overall sums valid domain memberships, so a multi-domain term can count more than once, and its mean and high-entropy percentage use that same denominator.}
\label{tab:entropy_distribution}
\end{table}

\subsection{Interval Reporting}\label{interval-reporting}

The statistics artifact reports per-term valid-entropy descriptives
(means and standard deviations, with exclusions counted) and the
inferential results in Table \ref{tab:pairwise_domain}; it does not
compute domain-level confidence intervals, so per-domain ambiguity-score
and context-variability intervals are not tabulated here. Exploratory
numerical differences are summarized by the Welch \(t\)-tests and the
omnibus ANOVA reported in Table \ref{tab:pairwise_domain}, with
per-domain entropy descriptives in Table \ref{tab:entropy_distribution}
and the summary in Figure \ref{fig:statistical_analysis}.

\subsection{Full-Text Parallel Layer}\label{full-text-parallel-layer}

A complementary descriptive analysis was run over 6997 open access full
texts harvested from PubMed Central (PMC), after excluding 76
correction, erratum and retraction notices and retracted items from 7073
stored records. The abstract corpus remains the headline corpus of this
study; the full-text layer is reported here as a separate convenience
sample with a higher extraction threshold. The analysis machinery is
shared: the same terminology extraction, domain assignment,
semantic-entropy, and CACE scoring implementations are applied to full
texts (Figure \ref{fig:fulltext_analysis}).

The layer comprises 44,476,595 tokens of running text, with a
per-document median of 5651 tokens. Table \ref{tab:fulltext_domain}
reports per-domain term counts and mean semantic entropy over the full
texts; 15 pairwise Welch \(t\)-tests (Benjamini-Hochberg corrected, as
in Table \ref{tab:pairwise_domain}) accompany the omnibus one-way ANOVA
on per-term semantic entropy, \(F = 7.7409\), \(p\)-value
\textless0.0001.

Anthropomorphic framing over the same full texts is scored as the
proportion of domain-term occurrence contexts containing an
anthropomorphic framing marker, using the same patterns as the abstract
layer: 0.0107 overall, with Economics at 0.0359 against 0.0074 for Unit
of Individuality and 0.0078 for Behavior \& Identity.

\begin{table}[h]
\centering
\begin{tabular}{|l|c|c|}
\hline
\textbf{Domain} & \textbf{Terms extracted} & \textbf{Mean $H$ (bits)} \\
\hline
Behavior \& Identity & 96 & 1.9700 \\
Economics & 35 & 1.7970 \\
Kin \& Relatedness & 24 & 1.9878 \\
Power \& Labor & 118 & 2.0102 \\
Sex \& Reproduction & 86 & 2.0381 \\
Unit of Individuality & 241 & 2.0274 \\
\hline
\end{tabular}
\caption{Per-domain terminology in the PMC full-text parallel layer: extracted-term counts and mean semantic entropy $H(t)$, computed over valid estimates with the same pipeline as the abstract layer; valid-estimate counts appear in Figure \ref{fig:fulltext_analysis}. Term extraction uses a higher minimum token frequency than the abstract layer because full texts are substantially longer.}
\label{tab:fulltext_domain}
\end{table}

\subsection{Discourse and Rhetorical
Layer}\label{discourse-and-rhetorical-layer}

A shared discourse stage counts rule-based markers of discourse
patterns, rhetorical strategies, argumentative structures, and
persuasive techniques in both layers. Texts shorter than 200 characters
are excluded. The abstract pass therefore covers 7583 texts (16 excluded
as too short; 100.0\% of the eligible texts). The full-text pass covers
1400 of 6997 texts, a deterministic 20.0\% sample of eligible texts
drawn to keep the pass computationally bounded. Figure
\ref{fig:discourse_comparison} compares the two layers; panels spanning
orders of magnitude use a symlog frequency axis.

Discourse patterns occur 1278 times as hierarchical framing in the
abstract layer against 1098 occurrences in the full-text layer, with
1045 versus 1088 economic metaphors, 32 versus 52 anthropomorphic
framings, and 27 versus 22 scale-ambiguous constructions.

Rhetorical-strategy marker counts in the two layers are: 1578 anecdotal
markers in the abstract layer versus 9446 in full texts, 246 versus
30,330 authority markers, 456 versus 1944 analogies, and 1148 versus
4204 generalizations.

The argumentative-structure pass identifies 1803 argumentative
structures in the abstract layer versus 1292 in the full-text sample.
Metaphorical-language markers occur 3830 times in the abstract layer
against 21,140 occurrences in full texts. All frequencies are raw marker
counts. The layers differ in size and sampling, so these counts are not
prevalence comparisons.

\clearpage

\subsection{Statistical and Source-Layer
Figures}\label{statistical-and-source-layer-figures}

Figures \ref{fig:statistical_analysis}--\ref{fig:arxiv_analysis}
summarize the headline statistics, the PMC full-text layer, the
cross-layer comparison, the discourse markers and the separate arXiv
preprint layer.

\begin{figure}[htbp]
\centering
\includegraphics[width=0.95\textwidth,height=0.78\textheight,keepaspectratio]{/Volumes/GitVault/Local/worktrees/ento-v1.3.1/ento_linguistics/output/figures/statistical_analysis.png}
\caption{Headline valid-entropy descriptives with nominal intervals, bias-corrected standardized differences, and exploratory ANOVA. Bar annotations report valid term counts; overlapping domain memberships and shared document contexts limit population inference.}
\label{fig:statistical_analysis}
\end{figure}

\begin{figure}[htbp]
\centering
\includegraphics[width=0.95\textwidth,height=0.78\textheight,keepaspectratio]{/Volumes/GitVault/Local/worktrees/ento-v1.3.1/ento_linguistics/output/figures/fulltext_analysis.png}
\caption{Separate PMC full-text statistics over all stored records by default. Extraction thresholds and context distributions differ from the abstract layer. Nominal intervals and multiplicity-adjusted threshold flags remain exploratory.}
\label{fig:fulltext_analysis}
\end{figure}

\begin{figure}[htbp]
\centering
\includegraphics[width=0.95\textwidth,height=0.78\textheight,keepaspectratio]{/Volumes/GitVault/Local/worktrees/ento-v1.3.1/ento_linguistics/output/figures/layer_comparison.png}
\caption{Mean valid context-cluster entropy in the abstract and PMC layers, with valid-term counts. This compares distinct convenience samples and extraction thresholds; it is not a matched robustness experiment or evidence of causal language effects.}
\label{fig:layer_comparison}
\end{figure}

\begin{figure}[htbp]
\centering
\includegraphics[width=0.95\textwidth,height=0.78\textheight,keepaspectratio]{/Volumes/GitVault/Local/worktrees/ento-v1.3.1/ento_linguistics/output/figures/discourse_comparison.png}
\caption{Lexical discourse, rhetorical and persuasive-pattern counts, with analyzed-text counts shown in the legend. PMC uses an explicitly bounded text sample. Raw frequencies depend on sample size and text length and must not be read as normalized prevalence or human-validated author intent.}
\label{fig:discourse_comparison}
\end{figure}

\begin{figure}[htbp]
\centering
\includegraphics[width=0.95\textwidth,height=0.78\textheight,keepaspectratio]{/Volumes/GitVault/Local/worktrees/ento-v1.3.1/ento_linguistics/output/figures/arxiv_analysis.png}
\caption{Separate arXiv preprint-layer descriptives and exploratory comparisons. Some groups have only one valid entropy estimate, for which intervals are omitted; very small groups and nearly zero within-group variance can yield large standardized differences without establishing generalizable effects.}
\label{fig:arxiv_analysis}
\end{figure}

\clearpage

\section{Supplemental Analysis: Theoretical
Extensions}\label{sec:supplemental_analysis}

These are proposed constructions, not additional experimental results.
The repository does not fit a generative model of scientific language,
measure variational free energy, estimate empirical Markov blankets, or
test a causal terminology intervention.

\subsection{Individuality and Conditional
Independence}\label{individuality-and-conditional-independence}

A Markov blanket specifies a conditional-independence relation between
internal variables \(\mu\) and external variables \(\eta\), conditional
on variables \(B\) \citep{friston2013life, kirchhoff2018markov}:

\begin{equation}\label{eq:markov_blanket}
P(\mu\mid\eta,B)=P(\mu\mid B).
\end{equation}

This relation is an assumption or a property to establish for a
specified probability model, on the support where the conditional
distributions are defined. A biological application requires named
variables, dynamics, observations, and tests of the proposed conditional
independence. A sensory/active decomposition requires further modeling
assumptions.

An ant's sensory and motor variables, or colony-level interaction
variables, could enter candidate models at different scales. Neither the
cuticle, a nest entrance, nor a pheromone field automatically
establishes a blanket. Likewise, calling a colony a superorganism is not
equivalent to proving this factorization. The Active Inferants
simulation provides related modeling context \citep{friedman2021active},
without validating a blanket inferred from terminology in this study.

\subsection{A Proposed Framing Score}\label{a-proposed-framing-score}

A future annotated study could define a bounded score:

\begin{equation}\label{eq:discursive_framing}
F_{\mathrm{frame}}(D,T)=\sum_{t\in T}w_t f_t(D)c_t,
\quad w_t\geq0,\quad\sum_{t\in T}w_t=1.
\end{equation}

Here \(f_t(D)\) and \(c_t\) would be explicitly defined annotation-based
quantities in \([0,1]\), giving a score in \([0,1]\) for nonempty \(T\).
Empty vocabularies would yield an undefined score, not an inferred
absence of framing. The weights require a prespecified rationale and
sensitivity analysis. This proposed quantity is distinct from
variational free energy and is not an output of the current
lexical-pattern pipeline.

\subsection{Ambiguity Annotation}\label{ambiguity-annotation}

Lexical, contextual, biological-scale, and temporal ambiguity are
candidate annotation categories. The implemented entropy measure
summarizes context-cluster occupancy; it does not independently assign
or validate these four categories. An annotation protocol should specify
meaning distinctions, uncertainty, annotator training, and agreement
measurement. Changes in context counts or source composition need
separate controls.

\subsection{Cross-Domain Mapping}\label{cross-domain-mapping}

Given an independently specified similarity score \(s(t,D_i,D_j)\), an
annotated cross-domain summary could be defined as:

\begin{equation}\label{eq:cross_domain_mapping}
M_{ij}=
\begin{cases}
|T_i\cap T_j|^{-1}\displaystyle\sum_{t\in T_i\cap T_j}s(t,D_i,D_j), & |T_i\cap T_j|>0,\\
\mathrm{undefined}, & |T_i\cap T_j|=0.
\end{cases}
\end{equation}

This is a proposed mean over shared terms. It is not the implemented
domain overlap coefficient, and an empty intersection does not establish
that two biological concepts are unrelated. Similarity and domain
membership require explicit definitions before interpretation.

\subsection{Temporal Networks}\label{temporal-networks}

A time-indexed analysis should first specify a common vertex registry,
matching vocabulary, and sampling rules. Let \(A(t)\) be weighted
adjacency matrices over that registry, with zero entries for absent
edges. Then an algebraically defined change is:

\begin{equation}\label{eq:temporal_network_evolution}
\Delta A(t)=A(t+1)-A(t).
\end{equation}

Changing matrix entries can reflect document availability, vocabulary
selection, source genre, or retrieval as well as changes in use.
Demonstrating a conceptual shift requires controls for these
alternatives and dated meaning annotation. The current BHL
literal-frequency analysis does not implement this longitudinal network
model.

\subsection{Multiscale Interpretation}\label{multiscale-interpretation}

Term, domain, and cross-domain summaries can organize different
observational scales. Claims about biological organization or scientific
thought require additional measurements at the corresponding scale.
Environment-Centric Active Inference and related interpretations remain
hypotheses requiring specified models and empirical tests, rather than
conclusions established by the present corpus.

\clearpage

\section{Supplemental Analysis: Case Studies and
Validation}\label{sec:supplemental_case_studies}

\subsection{Validation Agenda}\label{validation-agenda}

Expert classification review, interdisciplinary assessment, historical
interpretation, and cross-cultural comparison are proposed validation
activities. This repository does not report completed human-annotation
agreement, inter-rater reliability, multilingual validation, or a
prespecified subsampling and coefficient-sensitivity study. Software
tests and descriptive source-layer comparisons do not substitute for
those measurements.

\subsection{Historical Terminology Analysis}\label{sec:bhl_grounding}

The historical layer contains 2430 stored BHL mirror OCR documents dated
into three era buckets: 996 in 1850--1899, 1237 in 1900--1949, and 197
in 1950--1970. Retrieval selects volumes containing search expressions,
not a manually curated collection of ant-only works. Document titles,
dates, source identifiers, queries, and retrieval information are
recorded in the source sidecar.

\begin{table}[h]
\centering
\begin{tabular}{|l|r|r|r|}
\hline
\textbf{Term} & \textbf{1850--1899} & \textbf{1900--1949} & \textbf{1950--1970} \\
\hline
caste & 0.0735 & 0.1140 & 0.2253 \\
worker & 0.6394 & 0.7935 & 1.4169 \\
queen & 0.8861 & 1.0770 & 1.0064 \\
colony & 0.9283 & 1.1689 & 1.1047 \\
superorganism & 0.0003 & 0.0007 & 0.0051 \\
\hline
\end{tabular}
\caption{Complete-corpus literal OCR matches per 10,000 era tokens. Differences are descriptive of the selected volumes. OCR quality, topic, language, and spelling differences limit historical interpretation.}
\label{tab:superorganism_concept_evolution}
\end{table}

\begin{figure}[h]
\centering
\includegraphics[width=0.95\textwidth,height=0.78\textheight,keepaspectratio]{/Volumes/GitVault/Local/worktrees/ento-v1.3.1/ento_linguistics/output/figures/bhl_term_usage.png}
\caption{Full-corpus historical literal frequencies for six seed terms. Each panel uses its own rate axis. Neither zeros nor changing rates establish conceptual origin, prevalence in all entomological literature, or causal influence on research.}
\label{fig:bhl_term_usage}
\end{figure}

Rare matches for superorganism cannot establish that the concept
originated after 1970. Alternative spellings such as super-organism are
counted differently, and earlier theoretical discussion can exist
outside the sampled sources. Interpretation should be anchored in dated
sources such as \citet{wheeler1911}, rather than treating an OCR absence
as historical proof.

\subsection{Complete Document Coverage and Bounded
Entropy}\label{complete-document-coverage-and-bounded-entropy}

The BHL extraction/entropy/framing stack uses 996, 1237, and 197
complete documents, respectively. Default extraction and framing use all
era documents. Available and analyzed characters are recorded in the
artifact's coverage metadata. Entropy considers twenty frequent
extracted candidates per era using contexts across all documents; it is
a candidate-term sample rather than an exhaustive entropy census.
Candidate filtering can retain unrelated substring matches such as skin,
making, or queensland; unassigned candidates enter the overall sampled
entropy mean but not domain means. These are computational filter
outputs, not a curated historical vocabulary. An optional development
character budget produces a disclosed, separately fingerprinted document
subset. Neither default storage nor that subset establishes historical
representativeness.

Developmental evidence for canalized caste differentiation
\citep{qiu2022canalized} concerns biological mechanisms and does not
establish a historical linguistic trend. Evaluating that connection
would require a dated, annotated corpus with source-composition
controls.

\subsection{A Reproducible Case-Study
Protocol}\label{a-reproducible-case-study-protocol}

A subsequent study can select a specific term, reconcile each text to
its source, annotate meaning and biological referent in dated contexts,
and preregister comparisons between research traditions or eras.
Controls should distinguish genuine changes in use from retrieval, OCR,
spelling, and genre changes. CACE judgments should be gathered
independently of the automatic scoring vocabulary, with agreement and
uncertainty reported.

\clearpage

\section{Symbols and Notation Glossary}\label{sec:glossary}

This glossary defines the mathematical notation and domain-specific
terminology used throughout the manuscript.

\subsection{Mathematical Notation}\label{mathematical-notation}

{\def\LTcaptype{none} % do not increment counter
\begin{longtable}[]{@{}
  >{\raggedright\arraybackslash}p{(\linewidth - 4\tabcolsep) * \real{0.2500}}
  >{\raggedright\arraybackslash}p{(\linewidth - 4\tabcolsep) * \real{0.4062}}
  >{\raggedright\arraybackslash}p{(\linewidth - 4\tabcolsep) * \real{0.3438}}@{}}
\toprule\noalign{}
\begin{minipage}[b]{\linewidth}\raggedright
Symbol
\end{minipage} & \begin{minipage}[b]{\linewidth}\raggedright
Description
\end{minipage} & \begin{minipage}[b]{\linewidth}\raggedright
First Use
\end{minipage} \\
\midrule\noalign{}
\endhead
\bottomrule\noalign{}
\endlastfoot
\(G = (V, E)\) & Terminology network (graph with vertices and edges) &
Eq. \ref{eq:network_edge_weight} \\
\(D(t)\) & Set of Ento-Linguistic domains term \(t\) is assigned to &
Eq. \ref{eq:cace_appropriateness} \\
\(w(u,v)\) & Observed number of documents containing both terms \(u\)
and \(v\) & Eq. \ref{eq:network_edge_weight} \\
\(H(t)\) & Semantic entropy of term \(t\) in bits (Shannon entropy over
usage-context clusters) & Eq. \ref{eq:semantic_entropy} \\
\(H^*\) & Configured threshold (2.0 bits); not an independently
calibrated sense threshold & Eq. \ref{eq:semantic_entropy} \\
\(H_{\max}\) & Entropy ceiling for occupied clusters,
\(\log_2 k_{\mathrm{occupied}}\); zero for a single occupied cluster &
Eq. \ref{eq:semantic_entropy} \\
\(\hat{H}(t)\) & Normalized entropy \(H(t) / H_{\max}\) when
\(H_{\max}>0\); zero for one occupied cluster & Eq.
\ref{eq:semantic_entropy} \\
\(p_i\) & Empirical proportion of contexts assigned to semantic cluster
\(i\) & Eq. \ref{eq:semantic_entropy} \\
\(k\) & Requested number of computational clusters (\(k\)-means,
\(k = \max(2,\ \min(k_{\max},\ n{-}1,\ \max(2, \lfloor\!\sqrt{n}\rfloor)))\);
\(k_{\max}=5\), \(n=|C_t|\); \(k < n\)) & Eq.
\ref{eq:semantic_entropy} \\
\(C_t\) & Set of valid usage contexts of term \(t\) (sentences with more
than three words) & Eq. \ref{eq:semantic_entropy} \\
\(S_t\) & Set of biological scale levels expressed in term \(t\)'s
contexts & Eq. \ref{eq:cace_evolvability} \\
\(w_{AB}\) & Overlap coefficient (Szymkiewicz--Simpson) between concept
sets \(A\) and \(B\) & Eq. \ref{eq:overlap_coefficient} \\
\(\text{Clarity}(t)\) & CACE Clarity score:
\(\min(1,\max(0,1-H(t)/3.32))\) & Eq. \ref{eq:cace_clarity} \\
\(\text{Appropriateness}(t)\) & CACE Appropriateness score (penalizes
anthropomorphic terms) & Eq. \ref{eq:cace_appropriateness} \\
\(\text{Consistency}(t)\) & CACE Consistency score: mean pairwise cosine
similarity of context vectors & Eq. \ref{eq:cace_consistency} \\
\(\text{Evolvability}(t)\) & CACE Evolvability score: mean of heuristic
domain-breadth and scale-marker components & Eq.
\ref{eq:cace_evolvability} \\
\(\mathcal{A}\) & Set of anthropomorphic terms (queen, king, slave,
worker, soldier, nurse, \ldots) & Eq. \ref{eq:cace_appropriateness} \\
\(F_{\mathrm{frame}}(D, T)\) & Proposed bounded annotated framing score
for domain \(D\) and term set \(T\) & Supplemental Eq.
\ref{eq:discursive_framing} \\
\(M_{ij}\) & Cross-domain mapping strength between domains \(D_i\) and
\(D_j\) & Supplemental Eq. \ref{eq:cross_domain_mapping} \\
\(\Delta A(t)\) & Change in adjacency matrices on a common vertex
registry & Supplemental Eq. \ref{eq:temporal_network_evolution} \\
\(B\) & Markov Blanket boundary of a system & Supplemental Eq.
\ref{eq:markov_blanket} \\
\(\mu\) & Internal states (conditionally independent of external given
blanket) & Supplemental Eq. \ref{eq:markov_blanket} \\
\(\eta\) & External states & Supplemental Eq. \ref{eq:markov_blanket} \\
\end{longtable}
}

\subsection{Theoretical Terms}\label{theoretical-terms}

{\def\LTcaptype{none} % do not increment counter
\begin{longtable}[]{@{}
  >{\raggedright\arraybackslash}p{(\linewidth - 4\tabcolsep) * \real{0.2222}}
  >{\raggedright\arraybackslash}p{(\linewidth - 4\tabcolsep) * \real{0.4444}}
  >{\raggedright\arraybackslash}p{(\linewidth - 4\tabcolsep) * \real{0.3333}}@{}}
\toprule\noalign{}
\begin{minipage}[b]{\linewidth}\raggedright
Term
\end{minipage} & \begin{minipage}[b]{\linewidth}\raggedright
Definition
\end{minipage} & \begin{minipage}[b]{\linewidth}\raggedright
Context
\end{minipage} \\
\midrule\noalign{}
\endhead
\bottomrule\noalign{}
\endlastfoot
\textbf{Active Inference} & A corollary of the Free Energy Principle
stating that agents act to fulfill the predictions of their generative
models. & Sec. \ref{sec:introduction} \\
\textbf{CACE} & Clarity, Appropriateness, Consistency, Evolvability ---
four-dimensional meta-standard for evaluating scientific terminology. &
Sec. \ref{sec:methodology} \\
\textbf{Generative Model} & A probabilistic model of how sensory data is
generated from latent causes. & Sec. \ref{sec:discussion} \\
\textbf{Markov Blanket} & Variables rendering specified internal and
external variables conditionally independent in a probability model. &
Sec. \ref{sec:supplemental_analysis} \\
\textbf{Semantic Entropy} & Shannon entropy \(H(t)\) over the cluster
distribution of a term's usage contexts; describes computational cluster
occupancy without independently validated senses. & Sec.
\ref{sec:methodology} \\
\textbf{Superorganism} & A colony-level organismic concept; its
biological interpretation is not established by a term's frequency or by
an assumed blanket. & Sec. \ref{sec:introduction}; Sec.
\ref{sec:experimental_results} \\
\textbf{Variational Free Energy} & A variational bound on negative log
model evidence under a specified probabilistic model; not measured by
this corpus pipeline. & Sec. \ref{sec:discussion} \\
\end{longtable}
}

\subsection{Pipeline Modules}\label{pipeline-modules}

{\def\LTcaptype{none} % do not increment counter
\begin{longtable}[]{@{}
  >{\raggedright\arraybackslash}p{(\linewidth - 4\tabcolsep) * \real{0.3333}}
  >{\raggedright\arraybackslash}p{(\linewidth - 4\tabcolsep) * \real{0.3333}}
  >{\raggedright\arraybackslash}p{(\linewidth - 4\tabcolsep) * \real{0.3333}}@{}}
\toprule\noalign{}
\begin{minipage}[b]{\linewidth}\raggedright
Module
\end{minipage} & \begin{minipage}[b]{\linewidth}\raggedright
File
\end{minipage} & \begin{minipage}[b]{\linewidth}\raggedright
Function
\end{minipage} \\
\midrule\noalign{}
\endhead
\bottomrule\noalign{}
\endlastfoot
Text Processing & \texttt{text\_analysis.py} & Tokenization,
normalization, feature extraction \\
Term Extraction & \texttt{term\_extraction.py} & Domain-aware
terminology identification \\
Semantic Entropy & \texttt{semantic\_entropy.py} & Per-term \(H(t)\)
computation via TF-IDF + \(k\)-means \\
CACE Scoring & \texttt{cace\_scoring.py} & Four-dimensional terminology
evaluation \\
Domain Analysis & \texttt{domain\_analysis.py} & Per-domain framing and
ambiguity analysis \\
Conceptual Mapping & \texttt{conceptual\_mapping.py} & Cross-domain
concept graph construction \\
Rhetorical Analysis & \texttt{rhetorical\_analysis.py} & Framing
detection and argumentative scoring \\
Discourse Analysis & \texttt{discourse\_analysis.py} & Discourse pattern
classification \\
Statistics & \texttt{statistics.py} & Statistical validation
utilities \\
Visualization & \texttt{concept\_visualization.py} & Network and
domain-specific figure generation \\
\end{longtable}
}

\clearpage

\section{References}\label{sec:references}

\renewcommand{\bibsection}{}
\bibliography{references}

\clearpage

\end{document}
